% concatenated from main.tex
\documentclass{article}
\usepackage[utf8]{inputenc}

\usepackage[utf8]{inputenc}

\usepackage{comment}

\usepackage{graphicx} % Required for inserting images
\usepackage{amsmath,amsthm,verbatim,amssymb,amsfonts,amscd, graphicx}
\usepackage[utf8]{inputenc}
\usepackage{graphics}
\usepackage{hyperref}
\usepackage[baseline]{euflag}
\usepackage{enumitem}
\usepackage{apptools}
\usepackage{titlesec}
\usepackage{appendix}
\usepackage{xcolor}
\usepackage{dsfont}
\usepackage{authblk}

\AtAppendix{\counterwithin{lemma}{section}}

\topmargin0.0cm
\headheight0.0cm
\headsep0.0cm
\oddsidemargin0.0cm
\textheight23.0cm
\textwidth16.5cm
\footskip1.0cm
\theoremstyle{plain}
\newtheorem{prop}{}
\newtheorem{theorem}{Theorem}
\newtheorem{corollary}[theorem]{Corollary}
\newtheorem{lemma}[theorem]{Lemma}
\newtheorem{proposition}[theorem]{Proposition}
\newtheorem*{surfacecor}{Corollary 1}
\newtheorem{conjecture}[theorem]{Conjecture} 
\newtheorem{question}[theorem]{Question} 
\newtheorem{remark}[theorem]{Remark}
\newtheorem{example}[theorem]{Example} 

\theoremstyle{definition}
\newtheorem{definition}[theorem]{Definition}
\newcommand{\Ebig}[1]{%
    \E\bigl( #1\bigr)%
}
\newcommand{\Var}{\mathrm{Var}}
\newcommand{\RL}{\mathbb{R}}
\newcommand{\C}{{\mathcal{C}_{d}}}
\newcommand{\Cc} { \mathcal{C}(c_1,c_2,c_3)}
\newcommand{\IPn} {
I_{P_n}
}
\newcommand{\IPu} {
I_{P^u}
}
\newcommand{\inn}[2] {
\langle #1 , #2 \rangle
}
\newcommand{\Z}{\mathbb{Z}}
\newcommand{\bergstrom} {Bergstr{\"o}m }
\newcommand{\dd}{\,d}
\newcommand{\1}{\mathbf 1}


\usepackage[textwidth=2cm]{todonotes}

\newcommand\lampros[1]{\todo[size=\tiny, color=orange]{lg: #1}}

\DeclareMathOperator{\supp}{supp}



\title{Entropic analogues of Gr{\"u}nbaum's inequality}

\date{}

\author[1]{Matthieu Fradelizi}
\author[2]{Lampros Gavalakis}
\author[3]{Martin Rapaport}

\affil[1]{Univ Gustave Eiffel
                      }
                      
\affil[2]{University of Cambridge 
}

\affil[3]{
Carnegie Mellon University
                }
     


\begin{document}

\maketitle
 \footnotetext[1]{M.F. is with Univ Gustave Eiffel, Univ Paris Est Creteil, 
 CNRS, LAMA UMR8050 F-77447,
 Marne-la-Vall{\'e}e, France.\\
 email: matthieu.fradelizi@univ-eiffel.fr 
                      
L.G. is with the Department of Pure Mathematics and Mathematical Statistics, University of Cambridge,  UK and was supported by the Engineering and Physical Sciences Research Council [Grant Ref:EP/Y028732/1].\\
email: lg560@cam.ac.uk  

M.R. is with the 
Department of Mathematical Sciences, Carnegie Mellon University, Pittsburgh, 15213, PA, United States. \\
email: mrapapor@andrew.cmu.edu }
 
\begin{abstract} 
The classical Gr{\"u}nbaum inequality asserts that the proportion of the volume of a convex body cut off by a halfspace containing its barycenter is at least $1/e$. From its functional counterpart, for any log-concave random variable $X$, one has $\mathbb{P}(X\ge \mathbb{E}X)\ge 1/e$, with equality if and only if $X$ is exponential. 
Motivated by Gr{\"u}nbaum's inequality for convex bodies and its functional generalizations, we prove analogous inequalities for entropy, with characterizations of the equality cases. We show that if $X$ is a log-concave random variable on $\mathbb{R}$, then
$$
        h(X)-\frac{e}{e-1}H_2(1/e) \leq h(X|X \leq \mathbb{E}X) \leq h(X),
$$
where $h$ is the differential entropy, $H_2(\cdot)$ is the binary entropy function and $X|X\leq \mathbb{E}X$ stands for the distribution of $X$ conditional on $X\leq \mathbb{E}X$. We generalize the upper bound for all R{\'e}nyi entropies and the lower bound for {\em min}--entropy. Our inequalities are sharp and we characterize all equality cases. 

We discuss potential generalizations in high dimensions and give counterexamples in some directions. 

As an intermediate step for the proof of the lower bound, we establish a new inequality that we prove using a technique known as \textit{degrees of freedom}, combined with a standard KKT-type optimization lemma. Along the way, we characterize the equality case in a known comparison inequality between differential and {\em min}--entropy, which may be of independent interest.  
\end{abstract}

\section{ Introduction and main results}

Let $X \in \RL^n$ be a log-concave random vector. Let $H$ be an $(n-1)$-dimensional hyperplane through the mean of $X$, and denote by $H^-$ and $H^+$ the two half-spaces in $\mathbb{R}^n$ whose boundary is $H$. The functional version of Gr{\"u}nbaum's inequality \cite{grunbaum} states that
\begin{equation} \label{grunbaumeq}
    \frac{1}{e} \leq \mathbb{P}(X\in H^+),\hspace{0.2cm} \mathbb{P}(X\in H^-) \leq 1-\frac{1}{e}.
\end{equation}


In view of the analogy between entropy and volume \cite{costacover, dembocoverthomas, FMMZ}, the growing role of entropy methods in convex geometry \cite{ bobkovmadiman, bourgains,klartagstrong}, as well as the recently established entropic analogues of Bergstr{\"o}m's and Bonnesen's inequalities \cite{bergstromisit}, it is natural to ask whether Gr{\"u}nbaum's inequality admits an entropic counterpart. More precisely, we address the question of how the entropy of a log-concave distribution behaves when restricting to a half-space.

Let $X$ be a random vector with log-concave density $f$. For a hyperplane $H$, we write
\[
    X^+:=(X\mid X\in H^+)
    \quad \text{ and } \quad
    X^-:=(X\mid X\in H^-).
\]
Thus, for instance, $X^+$ has density
\[
    f^+(x)
    :=
    \frac{f(x)\mathbf 1_{H^+}(x)}
         {\mathbb P(X\in H^+)}.
\]
The differential entropy of $X$ is 
$$
h(X) = -\int_{\mathbb R^n}f(x) \log f(x) dx.
$$
We will abuse the notation and write $h(f) = h(X)$ when it is convenient. Denoting $p^+ := \mathbb{P}(X\in H^+)$ and $p^- := 1-p^+,$ a direct computation gives
\begin{equation} \label{elementaryid}
    h(X) = p^+h(X^+) +p^-h(X^-) + H_2(p^+), 
\end{equation}
where $H_2(p) := -p\log{p} - (1-p)\log{(1-p)}, p \in [0,1]$.
From an information theoretic standpoint, this can also be seen from the definition of mutual information (see \eqref{mutualinfodef} in the Appendix), since
\begin{align} \label{plusminuseq}
I(X;\mathbf{1}_{\{X\in H^+\}}) &= h(X)-h(X|\mathbf{1}_{\{X\in H^+\}}) = h(X) - p^+h(X^+) -p^-h(X^-) \\ \label{IH}
&= H(\mathbf{1}_{\{X\in H^+\}}) = H_2(p^+).
\end{align}

We aim to understand the relationship between $h(X^+), h(X^-)$ and $h(X)$. In the one-dimensional setting, we see from \eqref{elementaryid}, that if $H^+=\mathbb{R}_+$ and $f$ is symmetric around $0$, then 
$$
h(X^+) = h(X^-) = h(X)-\log{2}.
$$
On the other hand, for the centered exponential with density $f(x)=e^{-x-1}\mathbf{1}_{[-1,\infty)}(x),$ the memoryless property gives
\begin{equation} \label{exponentialvalue}
h(X^+) = h(X) = 1 = h(X^-) + \frac{e}{e-1}H_2\Bigl(\frac{1}{e}\Bigr),
\end{equation}
since this distribution also extremizes \eqref{grunbaumeq}.

Motivated by these examples, while thinking of the exponential distribution as an extreme point of the set of log-concave measures under moment constraints~(see, e.g., \cite{MNR}), it is natural to conjecture that in dimension $1$
\begin{equation} \label{condstrengthen}
    \max{\{h(X|X\in H^+),h(X|X\in H^-)\}} \leq h(X). 
\end{equation}

Our main contributions are twofold. First, we prove \eqref{condstrengthen} in dimension $1$, as well as the more general statement that
R\'enyi entropy of any order (defined just below) is non-increasing under one-sided truncation of a
one-dimensional log-concave density. Furthermore, we characterize all equality
cases. Second, under a centering assumption, we establish sharp reverse
inequalities for differential entropy and min-entropy, again with
complete equality characterizations.

We recall here that the R\'enyi entropy of order
$\alpha\in[0,\infty]$ of a random variable with density $f$ is
\begin{equation} \label{renyidef}
    h_\alpha(f)
    =
    \frac{1}{1-\alpha}
    \log\left(\int_{\mathbb R}f(x)^\alpha\,dx\right),
\end{equation}
with the usual limiting definitions at $\alpha=1$ and
$\alpha=\infty$. In particular,
\[
    h_1(f)=h(f)
    \quad \text{ and } \quad
    h_\infty(f)=-\log\|f\|_\infty,
\]
where $h_\infty$ is commonly referred to as min-entropy.

Our first main result is


%Since the exponential distribution is often an extreme point of the set of log-concave probability measures under moment constraints, as e.g. in \cite{MNR}, at least in dimension $1$ we also expect $\mathbf{1}_{[-1,\infty)}e^{-x-1}$ to maximize $h(X^+)$ and minimize $h(X^-)$. For the latter, we clearly need the zero-mean condition. 

%Indeed, we prove these claims in dimension $1$ in Theorem \ref{upperThIntro} below. 
%We recall here that the {\em R\'enyi} entropy of order $\alpha \in [0,\infty]$ is 
%$$
%h_{\alpha}(X) := \frac{1}{1-\alpha}\log{\Biggl( \int f(x)^{\alpha}\Biggr)},
%$$
%and $h_{\infty}(X) := -\log{\sup_{x\in \mathbb R}f(x)}$
%is commonly referred to as {\em min}--%entropy. 

\begin{theorem}
\label{upperThIntro}
Let $X \in \mathbb R$ have a log-concave density $f$. For $m \in \mathbb R$, let
\[
f^{-}_m(x):=\frac{f(x)\mathbf 1_{(-\infty,m]}(x)}{\int_{-\infty}^{m} f(y)\,dy}, \quad x \in \mathbb R,
\]
be the density of $X|X\le m$.
Then, for every $\alpha\in[0,\infty]$, and $m \in \mathbb R$ we have
\begin{equation} \label{alphaineqIntro}
h_\alpha\!\big(f^{-}_m\big)\le h_\alpha(f).
\end{equation}

Moreover, for any $\alpha \in (0,\infty]$, there is equality in \eqref{alphaineqIntro}, if and only if either the support of $X$ is contained in $(-\infty, m]$, or $X$ has an exponential distribution on $(-\infty,M]$ for some $M \geq m$. 
\end{theorem}

Repeatedly applying the conclusion of Theorem \ref{upperThIntro} shows that, for every $\alpha\in [0,+\infty]$, 
$$
    t \mapsto h_{\alpha}(X|X \leq t ) \text{ is non-decreasing in }t,
$$
or equivalently 
\begin{equation} \label{residualfunc}
    t \mapsto h_{\alpha}(X|X \geq t ) \text{ is non-increasing in }t.
\end{equation} 

The differential entropy case is the key step in the proof of
Theorem~\ref{upperThIntro}. The general R\'enyi inequality is deduced
from it by applying the differential entropy result to suitable normalized powers of the density. The equality analysis relies on a
rigidity result for the comparison inequality
$$
\log\|g\|_\infty \le 1+\int_{\mathbb R} g\log g,
$$
which was established in \cite[Proof of Theorem 7]{fradelizi2008increasing} and is an important ingredient in the proof of the inequality for differential entropy, Theorem \ref{entropygrunbaumTh}. We characterize all
equality cases in this inequality in
Proposition~\ref{prop:fradelizi-equality-1d}. As another consequence of
the min-entropy case of Theorem~\ref{upperThIntro} and Gr\"unbaum's
inequality, one recovers the known estimate
\cite[Theorem 4]{fradelizisections}, 
\[
    \|f\|_\infty\leq e f(0)
\]
for every centered log-concave probability density on $\mathbb R$.


In Section \ref{grunbaumineqsec} we break Theorem \ref{upperThIntro} into 
Proposition \ref{min-entroy}, 
Theorem \ref{entropygrunbaumTh},
Theorem \ref{thm:renyi-truncation}
and prove  
the equality case characterizations 
in 
Proposition 
\ref{cor:renyi-truncation-equality} using Proposition \ref{prop:fradelizi-equality-1d}. 
%In the latter, we characterize the equality case of the differential and {\em min}--entropy comparison inequality 
%$$
%\log\|g\|_\infty \le 1+\int_{\mathbb R} g\log g,
%$$
%which was established in \cite[Proof of Theorem 7]{fradelizi2008increasing} and is an important ingredient in the proof of the inequality for differential entropy, Theorem \ref{entropygrunbaumTh}.



%Another consequence is that, combining the {\em min}--entropy case of Theorem \ref{upperThIntro} with Gr{\"u}nbaum's inequality, we recover the following known comparison inequality
%{\color{red}To Do: Remove corollary}
%\begin{corollary} [\cite{fradelizisections}, Theorem 4] \label{matthieuCor}
%Let $f:\mathbb{R}\to[0,\infty)$ be a log-concave probability density with 
%\[
%\int_{\mathbb{R}} x\, f(x)\,dx = 0.
%\]
%ùThen
%\[
%\|f\|_\infty \le e\, f(0).
%\]
%\end{corollary}

  
Unlike the upper bound in Theorem~1, a reverse inequality requires the truncation point to be centered at the mean. Combining the $\alpha =1$ case of Theorem \ref{upperThIntro} (applied to $X^+$) with the identity \eqref{elementaryid} and Gr\"unbaum's
inequality gives the non-sharp reverse estimate
\[
        h(X)\le h(X^-) + e\,H_2(1/e).
\]
Our second main result identifies the sharp constant, both for
differential entropy and for min-entropy:

\begin{theorem}
\label{secondThIntro}
Let \(X\) be a random variable on $\mathbb R$ with zero mean and log-concave density \(f\), so that
\[
        \int_{\mathbb R} xf(x)\,dx=0.
\]
Let $\alpha \in \{1,\infty\}$ and $
f^{-}(x):=\frac{f(x)\mathbf 1_{(-\infty,0]}(x)}{\int_{-\infty}^{0} f(y)\,dy},  x \in \mathbb R
$.
Then
\begin{equation} \label{grunbaumineqintro}
        h_{\alpha}(f^-)
        \ge
        h_{\alpha}(f)-C_{\alpha},
\end{equation}
where 
$$
C_1 = \frac{e}{e-1}H_2(1/e)=\frac{e}{e-1}-\log(e-1) \quad \text{ and }  \quad C_{\infty} = \log(1+\sqrt{2}).
$$
For $\alpha = 1$, there is equality in \eqref{grunbaumineqintro} if and only if $X$ follows a centered exponential distribution on $[\frac{1}{\lambda}, \infty)$ for some $\lambda < 0$.

For $\alpha = \infty,$ there is equality in \eqref{grunbaumineqintro} if and only if $f$ is constant on a finite interval $[0,a], a \geq 0$ and exponential on $(-\infty, 0]$, or equivalently, if and only if $X^+$ is uniformly distributed on an interval and $X^-$ is exponentially distributed.
\end{theorem}
The extremizers in Theorem~\ref{secondThIntro} depend on the R\'enyi entropy
order. We remark first that the extremizer in \eqref{grunbaumineqintro} for $\alpha=1$ is the mirrored exponential to the one that extremizes \eqref{alphaineqIntro} in Theorem \ref{upperThIntro}, and  for $\alpha=+\infty$, the extremizer is instead
piecewise log-affine: it is exponential on one side of the origin and
constant on the other.


In Section \ref{grunbaumSec}, we break Theorem \ref{secondThIntro} into
Theorem \ref{minentropyreverseTh} for {\em min}--entropy and
Theorem \ref{thm:sharp-reverse-grunbaum-shannon} for differential entropy. 

The proof of the differential entropy inequality relies on a new
constrained entropy minimization result for log-concave densities on
$\mathbb R_+$. More precisely, an important ingredient is
Proposition~\ref{prop:normalized-variational-inequality}. This is reminiscent of the result of Melbourne, Nayar,
and Roberto \cite{MNR}, who identified one-sided exponentials as entropy
minimizers among log-concave random variables with fixed variance. Using the
\textit{degrees of freedom} method of \cite{fradeliziguedon}, we first reduce the
problem to densities that are log-affine on at most three intervals.
We then apply the Fritz--John stationarity conditions of
Lemma~\ref{FJ} to show that any minimizer must in fact be log-affine on
a single interval, for which the desired inequality can be verified
directly.

%An important ingredient in the proof for the differential entropy case is the inequality of  
%Proposition~\ref{prop:normalized-variational-inequality}. The degrees of freedom method of \cite{fradeliziguedon} allows us to reduce the inequality to a three-piece log-affine density. By applying the necessary stationarity conditions of Lemma \ref{FJ}, we deduce that any minimizer should in fact be one-piece log-affine, for which we verify the inequality directly. 

\medskip

{\bf Higher-dimensional extensions.} One might hope that Theorems \ref{upperThIntro} and \ref{secondThIntro} extend to higher dimensions. In particular, it is tempting to conjecture that \eqref{alphaineqIntro} and \eqref{grunbaumineqintro} remain true in any dimension. There are, however, two distinct obstructions.

First, Example \ref{counterexample} at the end of Section \ref{grunbaumineqsec}, shows that, already for min-entropy in dimension $2$, \eqref{alphaineqIntro} can fail if one does not require the boundary of the truncation to pass through the mean. 

Second, even under this additional assumption, the following CLT argument shows that \eqref{alphaineqIntro} cannot hold without an additive dimensional constant if the dimension is large enough. Indeed, let $X=(X_1,\ldots,X_n)$ where $X_1,\ldots ,X_n$   are independent centered exponential random variables, i.e., each $X_i$ has density
$
f_i(x)=e^{-(x+1)}\mathbf 1_{\{x\ge -1\}}(x).
$
Then \(X\) is centered and log-concave, and
\[
h(X)=n.
\]
Let $H^-=\{x=(x_1,\dots,x_n)\in\mathbb{R}^n; \sum_{i=1}^nx_i\le0\}$ and $X^-=X|X\in H^-.$
The density of \(X^-\)  is
\[
f_{X^-}(x)
=
\frac{e^{-\sum_{i=1}^n(x_i+1)}}
{\mathbb P(S_n\le 0)}
\mathbf 1_{\{x_i\ge -1,\ \sum_{i=1}^n{x_i}\le 0\}}.
\]
Therefore, letting $S_n:=\sum_{i=1}^n X_i$, we have
\[
h(X^-)
= -\int_{\mathbb R^n}f_{X^-}(x)\log{f_{X^-}(x)} \ dx =  
n+\mathbb E[S_n\mid S_n\le 0]
+
\log\mathbb P(S_n\le 0).
\]
By the CLT, as $n \to \infty,$
$
\mathbb P(S_n\le 0)\to \frac12,
$ 
and  
\[
\mathbb E[S_n\mid S_n\le0]
=-\sqrt{\frac{2n}{\pi}}+o(\sqrt n).
\]
Thus, 
\[
h(X^-)
=
n-\sqrt{\frac{2n}{\pi}}+o(\sqrt n)=h(X)-\sqrt{\frac{2n}{\pi}}+o(\sqrt n).
\]
Introducing $H^+=\{ x=(x_1,\dots,x_n)\in\mathbb{R}^n;\sum_{i=1}^nx_i\ge0\}$, $X^+=X|X\in H^+$, $p^{-}=\mathbb P(S_n\le 0)$ and $p^{+}=\mathbb P(S_n\ge 0)$. From \eqref{elementaryid}, we get
\[
h(X^+)=\frac{1}{p^+}\left(h(X)-p^{-}h(X^-)-H_2(p^+)\right)=h(X)+\sqrt{\frac{2n}{\pi}}+o(\sqrt n).
\]
%and hence
%\begin{equation} \label{Osqrtn}    h(X)-h(X\mid A)=O(\sqrt{n}).
%\end{equation}
Thus $h(X^+)-h(X)=\sqrt{\frac{2n}{\pi}}+o(\sqrt n)>0$ when $n$ is sufficiently large. Thus, even under the additional assumption that the truncating hyperplane contains the mean, there cannot be a direct dimensional extension of \eqref{alphaineqIntro} for $\alpha =1$.
%would then read $$h(X|A) \leq h(X).$$
%and by symmetry we would also have 
%$$h(X|A^{\mathrm C}) \leq h(X).$$
%But combining this with \eqref{elementaryid}, which holds in any dimension, would give 
%$$h(X) \leq h(X|A^{\mathrm{C}}) + \frac{1}{\mathbb P(A^{\mathrm{C}})}H_2(\mathbb P(A)) $$
%and therefore 
%$$h(X) = h(X|A^{\mathrm{C}}) + O(1) $$
%as $n \to \infty$, which contradicts \eqref{Osqrtn}. 
This leads to the following question.
\begin{question}
    What are the best constants $c_n, C_n>0$ such that for every log-concave random vector $X$ in $\mathbb R^n$ and for hyperplane $H$ containing $\mathbb E X$, one has
    $$
    h(X)-c_n\le h(X|X\in H^+) \leq h(X) + C_n ?
    $$
\end{question}
\noindent
We believe that the extremizers of the above question should be product exponential density as above. 
This leads to the following question, which we believe is of independent interest: 
\begin{question} \label{qtheta}
    Let $X = (X_1,\ldots,X_n)$ where $X_i$ are i.i.d. centered exponentials as above. Let $\theta \in \mathbb{S}^{n-1}$. Is it true that 
    $$
    \theta \mapsto h(X|\langle X,\theta \rangle \leq 0) 
    $$
    is minimized at the diagonal $\theta_1 := (1/\sqrt{n},\ldots,1/\sqrt{n})$ and maximized at the opposite diagonal $-\theta_1$?
\end{question}
\noindent
Question \ref{qtheta} is related in spirit to the treatment \cite{brazitikosexp}, as well as the results in \cite{eskenazisg,ENT} on the entropy of Gaussian mixtures. However, here we aim to extremize with respect to the conditioning event.   

The  argument above shows that any dimensional extension of Theorem \ref{secondThIntro} with constants $C_{\alpha,n}$ would have to satisfy $C_{1,n} = \Omega(\sqrt{n})$.
By letting $X$ be uniform on a centered convex body $K$, any such generalization of \eqref{grunbaumineqintro} would give a Gr{\"u}nbaum inequality of the form
$$
\mathrm{Vol(K^-)} \geq \frac{1}{e^{C_{\alpha,n}}}\mathrm{Vol}(K). 
$$
This estimate is suboptimal even in dimension $1$ (with the multiplicative constant $e^{-C_{\infty}}$) when compared to the optimal $\frac{1}{2}$ for an interval. Nevertheless, the constants in Theorem \ref{secondThIntro} are sharp, and we refer to \eqref{grunbaumineqintro} as an entropy analogue of Gr{\"u}nbaum's inequality (in dimension $1$). 
We do not know whether it is possible to recover the constant $\frac{1}{2}$ for intervals in dimension $1$ by a more clever application of the entropic estimate.

% It is worth noting that if we could generalize \eqref{grunbaumineqintro} to any dimension with the same constant, it would give an inferior constant $e^{-C_1} \approx 0.35$ when compared to Gr{\"u}nbaum's inequality, which gives $1/e \approx 0.368$. Nevertheless, the constant $C_1$ is tight for differential entropy in dimension $1$. 

\medskip
\textbf{Other connections and applications.} Combining Gr{\"u}nbaum's inequality \eqref{grunbaumeq} with the elementary property \eqref{IH}, we see that for a log-concave $X$, 
$$
I(X;\mathbf{1}_{\{X\in H^+\}}) \geq H_2\Bigl(\frac{1}{e}\Bigr).
$$
In a similar spirit, \eqref{condstrengthen} in $d=1$  strengthens the standard conditioning-reduces-entropy inequality (see \eqref{conditionreduceentropy} in the Appendix). 

Another way to state \eqref{grunbaumeq} from a statistical point of view is that trimming away a set of extreme values of probability at most $\frac{1}{e}$ will still keep the mean in the retained region, assuming the data distribution is log-concave. From this standpoint, \eqref{residualfunc} says that, in dimension one, assuming the data distribution is log-concave, removing a half-space of values never increases uncertainty. 



Let us remark that the function in \eqref{residualfunc} has been studied in the context of reliability engineering \cite{residual}. There, distributions with the property \eqref{residualfunc} are said to have decreasing uncertainty of residual life. Hence, our results imply that log-concave distributions in one dimension have decreasing uncertainty of residual life. The same function has been used as a basis for goodness-of-fit test for the exponential distribution~\cite{goodness}. See also \cite{hazard, piraeus, testinglifetime, chernoffresidual, propertiesofShannon, raocumulative} for this line of work. 




% \subsection{Bibliographical notes}
% \cite{propertiesofShannon} studies $h(X|t_1<X<t_2)$, says it is important for reliability engineering. Introduces order between distributions, defines a class for which this is non-increasing/decreasing. 
% \cite{raocumulative} introduces the cumulative residual entropy
% $-\int_{\mathbb{R}_+}\mathbb{P}(|X|>t)\log{\mathbb{P}(|X|>t)}dt$. 
% \cite{piraeus} recently introduces the cumulative residual interval entropy. 
% \cite{chernoffresidual} studies information measures of cumulative residual entropy. 


\subsection*{Statement on the use of A.I.}

The argument in the proof of Proposition \ref{prop:normalized-variational-inequality} and in particular the idea to use the stationarity conditions of Lemma \ref{FJ} to conclude that every local minimizer is one--piece log-affine was given by ChatGPT 5.4 and 5.5 after several prompts. The final proof was verified and written by the authors. 

\section{R{\'e}nyi entropies do not increase under truncation} \label{grunbaumineqsec}
%{\color{red} To Do: Explain structure, remove titles of Theorems. Put better section titles. Remove extra notation, $c_m, p, q, p_+,p_-$}

In this section, we prove Theorem~\ref{thm:renyi-truncation}.
We first establish the differential entropy case in Section \ref{diffentropysubsec}, from which the general
Rényi inequality follows. We prove the general inequality and then characterize the equality cases in Section \ref{renyisubsec}. We
conclude with a two-dimensional example illustrating the necessity
of the centering assumption in possible higher-dimensional extensions in Section \ref{counterexamplesec}.


\subsection{The differential entropy case} \label{diffentropysubsec}

\begin{lemma}
Let $f:\mathbb{R}\to[0,\infty)$ be a log-concave probability density. 
Then the function
\[
x\longmapsto -\frac{\int_{-\infty}^x f(y)\log f(y)\,dy}{\int_{-\infty}^x f(y)\,dy}+\log \Biggl(\int_{-\infty}^x f(y)\,dy\Biggr)
\]
is non-decreasing on $\mathbb{R}$.
\end{lemma}

\begin{proof}
Fix $x\in\mathbb{R}$ and define 
\[
F(x):=\int_{-\infty}^x f(y)\,dy,
\qquad
\Phi(x):=-\int_{-\infty}^x f(y)\log f(y)\,dy \quad \text{ and } \quad \psi(x):=\frac{\Phi(x)}{F(x)}.
\]
Consider also the probability density
\[
g_x(y):=\frac{f(y)}{F(x)}\mathbf{1}_{\{y\le x\}}.
\]
Since $f$ is log-concave, so is $g_x$. By the entropy bound for log-concave densities~\cite{fradelizi2008increasing}
\begin{equation} \label{matthieusineq}
\log\big(\|g_x\|_\infty\big)\le 1+\int g_x(y)\log g_x(y)\,dy.
\end{equation}
Noting that $\|g_x\|_\infty\ge g_x(x)=f(x)/F(x)$ and that
\[
\int g_x(y)\log g_x(y)\,dy
=
\frac{1}{F(x)}\int_{-\infty}^x f(y)\log f(y)\,dy-\log F(x),
\]
we obtain
\[
\log f(x)\le 1+\frac{1}{F(x)}\int_{-\infty}^x f(y)\log f(y)\,dy.
\]
A direct computation shows that
\[
\psi'(x)
=
\frac{f(x)}{F(x)}
\left(
-\log f(x)
+
\frac{1}{F(x)}\int_{-\infty}^x f(y)\log f(y)\,dy
\right),
\]
and the previous estimate yields
\[
\psi'(x)\ge -\frac{f(x)}{F(x)}.
\]
Therefore,
\[
(\psi+\log F)'(x)
=
\psi'(x)+\frac{f(x)}{F(x)}
\ge 0,
\]
which ends the proof.
\end{proof}

% \lampros{Change this to $\int_{-\infty}^xf(y)dy$ etc.}
% \begin{lemma}
% The function
% \[
% x\longmapsto \psi(x)+\log F(x)
% \]
% is non-increasing on $\mathbb{R}$. %In particular, for every $m\ge 0$,
% %\[
% %\psi(m)-\psi(0)\le \log F(0)-\log F(m).
% %\]
% \end{lemma}

% \begin{proof}


% Fix $x\in\mathbb{R}$ and define the probability density
% \[
% g_x(y):=\frac{f(y)}{F(x)}\mathbf{1}_{\{y\ge x\}}.
% \]
% Since $f$ is log-concave, so is $g_x$. By the entropy bound for log-concave densities due to Fradelizi
% (see []) 
% \[
% \log\big(\|g_x\|_\infty\big)\le 1+\int g_x(y)\log g_x(y)\,dy
% ,\]
% noting that $\|g_x\|_\infty\ge g_x(x)=f(x)/F(x)$ and that
% \[
% \int g_x(y)\log g_x(y)dy
% =
% \frac{1}{F(x)}\int_x^\infty f(y)\log f(y)\,dy-\log F(x),
% \]
% we obtain
% \[
% \log f(x)\le 1+\frac{1}{F(x)}\int_x^\infty f(y)\log f(y)\,dy.
% \]
% A direct computation shows that
% \[
% \psi'(x)
% =
% \frac{f(x)}{F(x)}
% \left(
% \log f(x)
% -
% \frac{1}{F(x)}\int_x^\infty f(y)\log f(y)\,dy
% \right),
% \]
% and the previous estimate yields $\psi'(x)\le f(x)/F(x)$. Combining the previous estimates, we conclude that
% \[
% (\psi+\log F)'(x)
% =
% \psi'(x)-\frac{f(x)}{F(x)}
% \le 0,
% \]
% which ends the proof.
% \end{proof}

%\textcolor{blue}{We say something %stronger in the introduction which is %true, the statement then should change}

%\textcolor{blue}{Mixtures in the integrals $m,0$ typos}
\begin{theorem}\label{entropygrunbaumTh}
Let $f$ be a log-concave probability density and for $m \in \RL$, let 
\[
f_{m}^-(x):=\frac{f(x)\mathbf 1_{(-\infty,m]}(x)}{\int_{-\infty}^{m} f(y)\,dy}.
\]
Then
\[
h(f_{m}^-)\le h(f).
\]
\end{theorem}

\begin{proof}  For \(x\in\mathbb R\), we define as before
\[
F(x):=\int_{-\infty}^x f(y)\,dy,\quad
\Phi(x):=-\int_{-\infty}^x f(y)\log f(y)\,dy \quad \text{ and } 
\psi(x):=\frac{\Phi(x)}{F(x)}.
\]
By the previous Lemma, the function
\[
x\mapsto \psi(x)+\log F(x)
\]
is non-decreasing. Denote $p_m^{-}:=\int_{-\infty}^{m} f(x)\,dx$. By definition, 
\[
\psi(m)+\log F(m)
=
-\frac1{p_m^{-}}\int_{-\infty}^m f(x)\log f(x)\,dx+\log p_m^{-}
=
h(f_{m}^-).
\]
On the other hand, since \(F(x)\to1\) as \(x\to\infty\), we have
\[
\lim_{x\to\infty}\big(\psi(x)+\log F(x)\big)
=
-\int_{\mathbb R} f(x)\log f(x)\,dx
=
h(f).
\]
Therefore,
\[
h(f_{m}^-)\le h(f).
\]


%For $x\ge 0$, define the left-tail functionals
%\[
%F_-(x):=\int_{-\infty}^{x} f(y)\,dy,
%\qquad
%\Phi_-(x):=-\int_{-\infty}^{x} f(y)\log f(y)\,dy,
%\qquad
%\psi_-(x):=\frac{\Phi_-(x)}{F_-(x)}.
%\]
%Applying the previous lemma to the reflected density, the function
%\[
%x\longmapsto \psi_-(x)+\log F_-(x)
%\]
%is non-decreasing on $[0,\infty)$. In particular,
%\[
%\psi_-(m)+\log F_-(m)\le \lim_{x\to\infty}\big(\psi_-(x)+\log F_-(x)\big).
%\]
%Since $F_-(m)=c_m$, we have
%\[
%\psi_-(m)+\log F_-(m)
%=
%-\frac{1}{c_m}\int_0^m f(x)\log f(x)\,dx+\log c_m
%=
%h(f^-).
%\]
%Moreover, as $x\to\infty$, $F_-(x)\to 1$ and
%\[
%\Phi_-(x)\to -\int_0^\infty f(x)\log f(x)\,dx,
%\]
%so that
%\[
%\lim_{x\to\infty}\big(\psi_-(x)+\log F_-(x)\big)=h(f).
%\]
%This concludes the proof.
\end{proof}

%\textcolor{blue}{eliminate all of these definitions since defined before...}

%For a real-valued continuous random variable $X$ with density $f$ on $\mathbb{R}$,
%the R\'enyi entropy of order $p \in [0,\infty]$ is defined by
%\[
%h_p(X)=h_p(f)
%=
%\frac{1}{1-p}
%\log\!\left(
%\int_{\mathbb{R}} f(x)^p\,dx
%\right),
%\]
%with the cases $p\in\{0,1,\infty\}$ understood in the usual limiting sense, namely
%\[
%h_0(f)=\log\operatorname{Leb}(\operatorname{supp} f),\qquad
%h_1(f)=-\int_{\mathbb{R}} f(x)\log f(x)\,dx,\qquad
%h_\infty(f)=-\log\|f\|_\infty.
%\]
\subsection{Rényi entropies and equality cases} \label{renyisubsec}

Since $\supp(f_{m}^-)\subset \supp(f)$, it is immediate that $h_{0}(f_{m}^-)\leq h_{0}(f)$. We first treat the min-entropy case.
\begin{proposition}%[Min-entropy decreases under truncation]
\label{min-entroy}
Let $f$ be a log-concave probability density and for $m \in \RL$, let 
\[
f_{m}^-(x):=\frac{f(x)\mathbf 1_{(-\infty,m]}(x)}{\int_{-\infty}^{m} f(y)\,dy}.
\]
Then
\[
h_{\infty}(f_m^-)\le h_{\infty}(f).
\]
Equivalently,
\[
\|f_m^-\|_\infty \ge \|f\|_\infty .
\]
\end{proposition}

\begin{proof}
Without loss of generality, assume that $m=0$ and define as shorthand
$
p_{-}:=\int_{-\infty}^{0} f(x)\,dx.$
We omit the index $0$ in $f_0^-$ and we write
\[
f^{-}(x):=\frac{f(x)\mathbf 1_{(-\infty,0]}(x)}{p_{-}}.
\]
By definition,
\[
\|f^-\|_\infty
=
\frac{1}{p_{-}}\sup_{x\le 0} f(x).
\]
Thus the claim is equivalent to
\[
\sup_{x\in\mathbb{R}} f(x)
\le
\frac{\sup_{x\le 0} f(x)}{\int_{-\infty}^{0} f(x)\,dx}.
\]

Let $t_{\max}\in\mathbb{R}$ be such that $f(t_{\max})=\sup_{x\in\mathbb{R}} f(x)$. If $t_{\max}\le 0$, then $\sup_{x\le 0} f(x)=\sup_{x\in\mathbb{R}} f(x)$ and $ p_{-}\le 1$, so
\[
\frac{\sup_{x\le 0} f(x)}{p_{-}}\ge \sup_{x\in\mathbb{R}} f(x),
\]
holds trivially. Assume now that $t_{\max}>0$.  
As before, define
\[
F(x):=\int_{-\infty}^{x} f(y)\,dy.
\]
Since $f$ is log-concave,  Pr\'ekopa's theorem implies that $F$ is log-concave on $\mathbb{R}$. Therefore $\log F$ is concave, and its derivative
\[
(\log F)'(x)=\frac{F'(x)}{F(x)}=\frac{f(x)}{F(x)}
\]
is non-increasing on $\mathbb{R}$. In particular, for $0\le t_{\max}$ we have
\[
\frac{f(0)}{F(0)}\ge \frac{f(t_{\max})}{F(t_{\max})}\geq \frac{\sup_{x\in \mathbb{R}}f(x)}{\int_{-\infty}^{\infty} f(x)dx}
=
\sup_{x\in\mathbb{R}} f(x).
\]
If there is equality then $\int_{-\infty}^{\infty} f(x)dx=F(t_{\max})$ and so $f(x) = 0$ on $(t_{\max},+\infty)$. Moreover $f/F=F'/F$ is constant on $[0,t_{\max}]$, thus there exists $c\in\mathbb{R}$ such that $F(t)=F(0)e^{ct}$, for every $t\in[0,t_{\max}]$. Hence $f(t)=F'(t)=F(0)ce^{ct}$. By log-concavity of $f$, this implies that $f(t)\le F(0)ce^{ct}$, for every $t\le t_{\max}$. Hence, integrating, we deduce that $F(t)\le F(0)e^{ct}$, for every $t\le t_{\max}$. Since there is equality for $t=t_{\max}$, we conclude that $f(t)= F(0)ce^{ct}$, for every $t\le t_{\max}$.
\end{proof}

As an immediate consequence of Proposition \ref{min-entroy} and the functional
form of Gr\"unbaum's inequality, we recover the estimate
\[
\|f\|_\infty \leq e f(0)
\]
for every centered log-concave probability density $f$ on $\mathbb R$;
see \cite[Theorem~4]{fradelizisections}.
Indeed, without loss of generality, assume that a maximizer $t_{\max}$ of $f$ satisfies $t_{\max}\ge 0$.
By Proposition~\ref{min-entroy}, we have
\[
\sup_{x\in\mathbb{R}} f(x)\le \frac{f(0)}{F(0)},
\qquad
F(0)=\int_{-\infty}^{0} f(x)\,dx.
\]

The functional Grünbaum inequality (see \cite[Lemma 5.12]{lovasz}) states that for any log-concave probability density on
$\mathbb{R}^n$ and any half-space $H$ containing its centroid,
\[
\int_H f(x)\,dx \ge \frac{1}{e}.
\]
Applying this inequality in dimension $n=1$ with centroid $0$ and halfspace $H=(-\infty,0]$, we obtain
$F(0) \ge \frac{1}{e}.$ Thus,
\[
\sup_{x\in\mathbb{R}} f(x)\le \frac{f(0)}{F(0)}\le e\,f(0).
\]



\medskip

Recall the definition of R{\'e}nyi entropy in \eqref{renyidef}.
We have treated the endpoint cases ($\alpha=0,\alpha=1$ and $\alpha=\infty$). Next we state the general result: 
\begin{theorem}%[R\'enyi entropy decreases under truncation]
\label{thm:renyi-truncation}
Let $f$ be a log-concave probability density and for $m \in \RL$, let 
\[
f_{m}^-(x):=\frac{f(x)\mathbf 1_{(-\infty,m]}(x)}{\int_{-\infty}^{m} f(y)\,dy}.
\]
Then, for every $\alpha\in[0,\infty]$, 
\[
h_\alpha\!\big(f_{m}^{-}\big)\le h_\alpha(f).
\]
\end{theorem}

\begin{proof}
The cases $\alpha=0$ and $\alpha=\infty$ were established above.
We may therefore assume that $\alpha\in(0,\infty)$ and, without loss of generality, take $m=0$.
We denote $\mathbb R^- := (-\infty,0]$.  
Define
\[
f^-(x) := \frac{f(x)\mathbf 1_{\mathbb R^-}(x)}{p^{-}},
\qquad
p^{-} := \int_{\mathbb R^-} f(x)\,dx,
\]
and
for $\alpha \ge 0$, let
\[
G(\alpha)
:= \frac{\|f\|_\alpha}{\|f\,\mathbf 1_{\mathbb R^-}\|_\alpha}
= \left(
\frac{\int_{\mathbb R} f(x)^\alpha\,dx}{\int_{\mathbb R^-} f(x)^\alpha\,dx}
\right)^{1/\alpha}.
\]

\medskip

We claim that the function $\alpha \mapsto G(\alpha)$ is non-increasing on $(0,\infty)$.
It suffices to show that $\frac{d}{d\alpha}\log G(\alpha)\le 0$.  
We first write
\[
\log G(\alpha)
=
\frac{1}{\alpha}\log\!\int_{\mathbb R} f(x)^\alpha\,dx
-
\frac{1}{\alpha}\log\!\int_{\mathbb R^-} f(x)^\alpha\,dx.
\]
A direct differentiation yields
\begin{equation}\label{derivada}
\frac{d}{d\alpha}\log G(\alpha)
=
\frac{1}{\alpha^2}
\Bigg[
\alpha\,\frac{\int_{\mathbb R} f(x)^\alpha \log f(x)\,dx}{\int_{\mathbb R} f(x)^\alpha\,dx}
-
\log\!\int_{\mathbb R} f(x)^\alpha\,dx
-
\alpha\,\frac{\int_{\mathbb R^-} f(x)^\alpha \log f(x)\,dx}{\int_{\mathbb R^-} f(x)^\alpha\,dx}
+
\log\!\int_{\mathbb R^-} f(x)^\alpha\,dx
\Bigg].
\end{equation}
To interpret \eqref{derivada} in terms of differential entropies, define the probability densities
\begin{equation} \label{powered}
g_\alpha(x)
:=
\frac{f(x)^\alpha}{\int_{\mathbb R} f(y)^\alpha\,dy},
\qquad
g_\alpha^-(x)
:=
\frac{f(x)^\alpha \mathbf 1_{\mathbb R^-}(x)}{\int_{\mathbb R^-} f(y)^\alpha\,dy}.
\end{equation}
A straightforward computation shows that
\[
h(g_\alpha)
=
-\int_{\mathbb R} g_\alpha(x)\log g_\alpha(x)\,dx
=
-\alpha\,\frac{\int_{\mathbb R} f(x)^\alpha \log f(x)\,dx}{\int_{\mathbb R} f(x)^\alpha\,dx}
+
\log\!\int_{\mathbb R} f(x)^\alpha\,dx,
\]
and similarly,
\[
h(g_\alpha^-)
=
-\int_{\mathbb R^-} g_\alpha^-(x)\log g_\alpha^-(x)\,dx
=
-\alpha\,\frac{\int_{\mathbb R^-} f(x)^\alpha \log f(x)\,dx}{\int_{\mathbb R^-} f(x)^\alpha\,dx}
+
\log\!\int_{\mathbb R^-} f(x)^\alpha\,dx.
\]
Substituting into \eqref{derivada} yields the identity
\begin{equation} \label{Gder}
\frac{d}{d\alpha}\log G(\alpha)
=
-\frac{1}{\alpha^2}
\big(
h(g_\alpha)-h(g_\alpha^-)
\big).
\end{equation}

Since $f$ is log-concave, $f^\alpha$ is log-concave for all $\alpha> 0$, and hence $g_\alpha$ is a log-concave probability density on $\mathbb R$.  
Applying Theorem \ref{entropygrunbaumTh} to $g_\alpha$ and its restriction $g_\alpha^-$, we obtain
\[
h(g_\alpha^-)\le h(g_\alpha).
\]
Together with \eqref{Gder}, this implies
\[
\frac{d}{d\alpha}\log G(\alpha)\le 0,
\]
so that $G(\alpha)$ is non-increasing on $(0,\infty)$.
Consequently,
\[
G(\alpha)\le G(1)\qquad\text{for all }\alpha\ge 1,
\]
and 
\[
G(\alpha)\geq G(1)\qquad\text{for all }\alpha\leq 1,
\]
which is in both cases equivalent to
\[
h_\alpha(f^-)\le h_\alpha(f).
\]
%The case $\alpha = 0$ is straightforward. 
\end{proof}




%\subsection{Equality cases }

We now turn to the equality cases. We will need the following lemma:
\begin{lemma}[\cite{santalo}, Lemma 3]
\label{lem:affine-rigidity}
Let $\phi:\mathbb{R}\to \mathbb{R}\cup\{+\infty\}$ be a convex function such that $
\int_{\mathbb{R}} e^{-\phi(t)}\,dt < +\infty.$
Let
\[
U:=\{\phi<+\infty\}.
\]
Suppose that for some $x\in U$ and for almost every $y\in U$, one has
\[
\phi(x)=\phi(y)+(x-y)\phi'(y).
\]
Then for every $z\in U$, the function $\phi$ is affine on the segment $[x,z]$.
\end{lemma}

The next proposition gives the one-dimensional equality cases in 
\begin{equation} \label{matthieu}
\log\|g\|_\infty \le 1+\int_{\mathbb R} g\log g,
\end{equation}
which is established in \cite[Proof of Theorem 7]{fradelizi2008increasing}.




\begin{proposition}
\label{prop:fradelizi-equality-1d}
Let $g:\mathbb R\to[0,\infty)$ be a log-concave probability density, and write
$g=e^{-\phi},$ where $\phi:\mathbb R\to \mathbb R\cup\{+\infty\}$ is convex.
Then
\[
\log\|g\|_\infty = 1+\int_{\mathbb R} g\log g
\]
if and only if

\begin{equation} \label{formula}
g(t)=\frac{1}{\alpha+\beta}
\begin{cases}
e^{\frac{t-t_0}{\alpha}}, & t\le t_0,\\[1mm]
e^{-\frac{t-t_0}{\beta}}, & t\ge t_0,
\end{cases}
\qquad\text{for some }t_0 \in \mathbb R, (\alpha, \beta) \in \mathbb R_+^2 \setminus \{(0,0)\}.
\end{equation}
\end{proposition}
Before giving the proof, we note that \eqref{formula} includes the cases where $\alpha$ or $\beta$ are equal to $0$ (but not both). In other words, the equality cases are precisely the two one-sided exponential
families and the asymmetric double-exponential family.
\begin{proof}
Let $t_0$ be a point where $g$ attains its maximum, equivalently where $\phi$
attains its minimum.
Since $\|g\|_\infty=e^{-\phi(t_0)}$ and $
\int_{\mathbb R} g\log g
=
-\int_{\mathbb R}\phi(y)e^{-\phi(y)}\,dy,$ 
the equality
\[
\log\|g\|_\infty = 1+\int g\log g
\]
is equivalent to
\begin{equation}\label{eq:frad-equality}
\phi(t_0)=\int_{\mathbb R}\phi(y)e^{-\phi(y)}\,dy-1.
\end{equation}
Let
\[
U:=\{\phi<+\infty\},
\qquad
a:=\inf U,\qquad b:=\sup U.
\]
We proceed as in the proof of \cite[Theorem 7]{fradelizi2008increasing}. Convexity gives, for every $y\in U$,
\[
\phi(t_0)\ge \phi(y)+(t_0-y)\phi'(y).
\]
Multiplying by $e^{-\phi(y)}$ and integrating over $U$, we obtain
\begin{equation}\label{eq:frad-step1}
\phi(t_0)
\ge
\int_U \phi(y)e^{-\phi(y)}\,dy
+
\int_U (t_0-y)\phi'(y)e^{-\phi(y)}\,dy.
\end{equation}
An integration by parts yields
\[
\int_U (t_0-y)\phi'(y)e^{-\phi(y)}\,dy
=
\Big[-(t_0-y)e^{-\phi(y)}\Big]_{a}^{b}
-\int_U e^{-\phi(y)}\,dy.
\]
Since $\int_U e^{-\phi}=1$, this becomes
\begin{equation}\label{eq:frad-step2}
\int_U (t_0-y)\phi'(y)e^{-\phi(y)}\,dy
=
-(t_0-b)e^{-\phi(b)}+(t_0-a)e^{-\phi(a)}-1
\ge -1.
\end{equation}
Combining \eqref{eq:frad-step1} and \eqref{eq:frad-step2} gives
\[
\phi(t_0)\ge \int_U \phi(y)e^{-\phi(y)}\,dy-1
=
\int_{\mathbb R}\phi(y)e^{-\phi(y)}\,dy-1.
\]
This is precisely the inequality
\[
\log\|g\|_\infty \le 1+\int_{\mathbb R}g\log g.
\]
Moreover, equality in this inequality is equivalent to \eqref{eq:frad-equality}.
Assume now that equality holds in \eqref{eq:frad-equality}. Then equality must
hold both in \eqref{eq:frad-step1} and \eqref{eq:frad-step2}.
Equality in
\eqref{eq:frad-step1} implies
\[
\phi(t_0)=\phi(y)+(t_0-y)\phi'(y)
\qquad\text{for almost every }y\in U,
\]
because the integrand
\[
\phi(t_0)-\phi(y)-(t_0-y)\phi'(y)
\]
is nonnegative. By Lemma~\ref{lem:affine-rigidity}, $\phi$ is affine on every
segment $[t_0,z]\subset U$. In dimension one this means that $\phi$ is affine
on $[a,t_0]$ and on $[t_0,b]$. Equality in \eqref{eq:frad-step2} means precisely that the boundary term vanishes:
\begin{equation}\label{eq:boundary-vanish}
(t_0-a)e^{-\phi(a)}=0,
\qquad
(b-t_0)e^{-\phi(b)}=0.
\end{equation}
There are exactly two regimes:

\smallskip

\noindent\emph{Case (i):} Suppose first that $t_0$ coincides with one endpoint of the support. If
$t_0=a$, then the support is of the form $[t_0,\infty)$, and since $\phi$ is
affine on the whole support,
\[
\phi(t)=\phi(t_0)+\lambda(t-t_0)
\qquad (t\ge t_0)
\]
for some $\lambda>0$. Hence
\[
g(t)=\lambda e^{-\lambda(t-t_0)}\mathbf 1_{[t_0,\infty)}(t).
\]
This corresponds to the displayed formula with $\alpha=0$ and
$\beta=1/\lambda$. The case $t_0=b$ is symmetric and gives
\[
g(t)=\lambda e^{\lambda(t-t_0)}\mathbf 1_{(-\infty,t_0]}(t),
\]
which corresponds to the displayed formula with $\alpha=1/\lambda$ and
$\beta=0$.

\smallskip

\noindent\emph{Case (ii):} Suppose now that $t_0$ does not coincide with an endpoint of the support. Then
\eqref{eq:boundary-vanish} forces
\[
e^{-\phi(a)}=e^{-\phi(b)}=0.
\]
Since $\phi$ is affine on each side of $t_0$, this implies that $a=-\infty$ and
$b=+\infty$. Therefore
\[
\phi(t)=\phi(t_0)+\frac{t_0-t}{\alpha}
\quad (t\le t_0),
\qquad
\phi(t)=\phi(t_0)+\frac{t-t_0}{\beta}
\quad (t\ge t_0)
\]
for some $\alpha,\beta>0$. Hence
\[
g(t)=C
\begin{cases}
e^{\frac{t-t_0}{\alpha}}, & t\le t_0,\\[1mm]
e^{-\frac{t-t_0}{\beta}}, & t\ge t_0,
\end{cases}
\]
and normalization gives
\[
C=\frac{1}{\alpha+\beta}.
\]

Conversely, a direct computation shows that every density described above
satisfies
\[
\log\|g\|_\infty = 1+\int g\log g .
\]
This completes the proof.
\end{proof}

% We now identify the equality cases in Theorem~\ref{entropygrunbaumTh}.

% \begin{proposition}[Equality case in Theorem~\ref{entropygrunbaumTh}]
% \label{prop:equality-entropy-restriction}
% Let $f:\mathbb R\to[0,\infty)$ be a log-concave probability density, and let
% \[
% c_m:=\int_{-\infty}^m f(x)\,dx,
% \qquad
% f^-(x):=\frac{f(x)\mathbf 1_{(-\infty,m]}(x)}{c_m}.
% \]
% Then
% \[
% h(f^-)=h(f)
% \]
% if and only if one of the following holds:
% \begin{enumerate}
% \item[(a)] $c_m=1$, equivalently $\operatorname{supp}(f)\subset(-\infty,m]$;
% \item[(b)] there exist $\lambda>0$ and $d\ge m$ such that
% \[
% f(x)=\lambda e^{\lambda(x-d)}\mathbf 1_{(-\infty,d]}(x).
% \]
% \end{enumerate}
% \end{proposition}

% \begin{proof}

% \end{proof}
Finally we characterize the equality cases in Theorem \ref{thm:renyi-truncation}]:
\begin{proposition}
\label{cor:renyi-truncation-equality}
Let $f$ be a log-concave probability density and for $m \in \RL$, let 
\[
f_{m}^-(x):=\frac{f(x)\mathbf 1_{(-\infty,m]}(x)}{\int_{-\infty}^{m} f(y)\,dy}.
\]
Let also $\alpha\in(0,\infty]$. Then for any $m \in \RL$,
\[
h_\alpha(f_m^-)=h_\alpha(f)
\]
if and only if either $\int_{-\infty}^{m} f(y)\,dy =1$, or there exist $\lambda>0$ and $d\ge m$ such that
\[
f(x)=\lambda e^{\lambda(x-d)}\mathbf 1_{(-\infty,d]}(x).
\]
\end{proposition}

\begin{proof}
Consider first the differential entropy case $\alpha =1$ and define, as before, $p_m^- := \int_{-\infty}^{m} f(y)\,dy$. If $p_m^{-}=1$, then $f_m^-=f$ a.e., so equality is immediate. Assume now that $p_m^{-}<1$ and that
\[
h(f_m^-)=h(f).
\]
Let
\[
d:=\sup\operatorname{supp}(f).
\]
For $x\in[m,d)$, define
\[
F(x):=\int_{-\infty}^x f(y)\,dy,
\qquad
g_x(y):=\frac{f(y)\mathbf 1_{(-\infty,x]}(y)}{F(x)}.
\]
By the proof of Theorem~\ref{entropygrunbaumTh}, the equality $h(f_m^-)=h(f)$ forces equality in   
inequality~\eqref{matthieusineq}
for every truncated density $g_x$, $x\in[m,d)$. 
By Proposition~\ref{prop:fradelizi-equality-1d}, each $g_x$ must belong to one of the equality regimes described there. Since
\[
\operatorname{supp}(g_x)\subset(-\infty,x],
\]
the asymmetric double-exponential regime is impossible, and the only possible
one-sided regime is the exponential supported on a left half-line. Hence, for every $x\in[m,d)$, there exists $\lambda_x>0$ such that
\[
g_x(y)=\lambda_x e^{\lambda_x(y-x)}\mathbf 1_{(-\infty,x]}(y).
\]
Equivalently,
\[
f(y)=F(x)\lambda_x e^{\lambda_x(y-x)}
\qquad (y\le x).
\]
Comparing these identities for different values of $x$, we conclude that $\lambda_x$ does not depend on $x$. It follows that $f$ must be a one-sided exponential supported on a left half-line, namely
\[
f(x)=\lambda e^{\lambda(x-d)}\mathbf 1_{(-\infty,d]}(x)
\]
for some $\lambda>0$ and some $d\ge m$. The converse is immediate by the memoryless property of the exponential density.

Now let $0<\alpha \neq 1$. By translation invariance, we may assume $m=0$. Let
\[
p^{-}:=\int_{-\infty}^0 f(x)\,dx,
\qquad
f^-(x):=\frac{f(x)\mathbf 1_{(-\infty,0]}(x)}{p^{-}}
\quad \text{and} \quad  
G(\alpha):=\frac{\|f\|_\alpha}{\|f\mathbf 1_{(-\infty,0]}\|_\alpha}.
\]


Suppose $h_\alpha(f^-)=h_\alpha(f)$. 
As in the proof of
Theorem~\ref{thm:renyi-truncation}, we deduce that $G(\alpha)=G(1)$. Since
$G$ is non-increasing on $(0,\infty)$, the function $G$
must therefore be constant on $[\min\{\alpha,1\},\max\{\alpha,1\}]$. By \eqref{Gder}, we have equality in $h(g_{\alpha}^-) = h(g_{\alpha})$, where $g_{\alpha}$ is defined as in \eqref{powered}. By the differential entropy case, we conclude that
either $\operatorname{supp}(g_{\alpha})\subset(-\infty,0]$, i.e.\ $p^{-}=1$, or else $g_{\alpha}$ is a one-sided exponential density. Since $g_\alpha\propto f^\alpha$, the latter implies that
\[
f(x)=\lambda e^{\lambda(x-d)}\mathbf 1_{(-\infty,d]}(x)
\]
for some $\lambda>0$ and $d\ge0$. The converse is again immediate by the memoryless property of the exponential density.  The case \(\alpha=\infty\) was already proved in Proposition \ref{min-entroy}.
\end{proof}

\subsection{The centering assumption in higher dimensions}
\label{counterexamplesec}

Finally, we record the example mentioned in the Introduction, which shows that, already for min-entropy in dimension
two, one cannot expect an analogous statement without a centering assumption on the
cutting hyperplane.
%\textcolor{blue}{Perhaps move it somewhere else?... now one %could believe that this example is an example of Theorem 7}

\begin{example} \label{counterexample}
Let \(a\geq1\), and set
\[
H=\{x\in\mathbb R^2:x_1+x_2=1\},\qquad
H^+=\{x\in\mathbb R^2:x_1+x_2\ge 1\},
\qquad
H^-=\{x\in\mathbb R^2:x_1+x_2\le 1\}.
\]
Consider $f(x) = (1-\frac{x_1}{a}-\frac{x_2}{a})_+, x_1,x_2 \geq 0.$ Then 
$$
\int_{\mathbb{R}^2} f = \frac{a^2}{6},\hspace{0.2cm} \max_{x \in H^+}f(x) = 1-\frac{1}{a}.
$$
On the other hand, direct computation shows 
$$
\int_{H^-}f = \frac{1}{2}-\frac{1}{3a}
$$
and thus
$$
\int_{H^+}f = \frac{a^2}{6}-\frac12+\frac1{3a}.
$$
The inequality 
\begin{equation} \label{ineqinexample}
\frac{\max_{H^+}f}{\max f} \geq \frac{\int_{H^+}f}{\int_{\mathbb{R}^2} f}
\end{equation}
is equivalent to 
$$
1 \leq \frac3{a}-\frac2{a^2}
$$
which is 
$$
(a-1)(a-2) \leq 0,
$$
i.e. satisfied for $a \in [1,2].$
One verifies that
$$
\frac{\int_{\mathbb{R}^2} xf(x)dx}{\int_{\mathbb{R}^2} f(x)dx} = (a/4,a/4).
$$

Hence $\int \frac{xf(x)}{\int f} dx \in H$ if and only if $a=2$. Consequently, taking any \(a>2\), the barycenter of \(f\) does not lie on \(H\), and the inequality
\eqref{ineqinexample} fails. This shows that, in dimension two, the condition that the cutting hyperplane pass through the barycenter cannot simply be omitted in the min-entropy truncation inequality.


%Since, $\int \frac{xf(x)}{\int f} dx \in H$ is equivalent to $(a/4,a/4) \in \{x_1+x_2=1\}$, which holds if and only if $a=2$, we see that the inequality \eqref{ineqinexample} can fail in general if the barycenter of $f$ does not lie in $H$.
\end{example}


\section{Proof of Gr\"unbaum inequalities for differential and min--entropy}
\label{grunbaumSec}
%\textcolor{blue}{Sharp reverse Grünbaum inequalities?}

Let $f$ be a zero-mean log-concave probability density on $\mathbb R$,
and 
\[
        f^+(x):=
        \frac{f(x)\mathbf 1_{\{x\ge 0\}}}{\int_0^\infty f(y)\,dy}.
\]
In this section, we would like to understand, for every $\alpha>0$, what is the best 
constant $C_\alpha<\infty$ such that
\[
        h_\alpha(f^+) \ge h_\alpha(f)-C_\alpha .
\]
Such a constant exists by monotonicity of R{\'e}nyi entropy together with the reverse comparison for log-concave densities~\cite{fradelizilimadiman} and Theorem \ref{thm:sharp-reverse-grunbaum-shannon} for differential entropy that we prove later in this section. 

Since $h_0(f)=\log |\operatorname{supp} f|$, no uniform finite constant should
hold at $\alpha=0$. Thus, we must have
\[
        C_\alpha\longrightarrow +\infty
        \qquad\text{as }\alpha\downarrow 0.
\]


Moreover, as we will see later in this section, the extremizers are different for $\mathrm{min}$--entropy and for differential entropy. This indicates that a unified reverse result is more subtle to obtain and we expect the extremizers to be some interpolating family of densities. Below we establish sharp reverse inequalities for the limiting cases $\alpha=1,\infty$. We will need the following lemmas.


\begin{lemma} [Lemma of Section 2.1. in \cite{milmanisotropic}]
\label{lem:moment-lower-decreasing}
Let $g:[0,\infty)\to[0,\infty)$ be integrable. 
\[
        \int_0^\infty t g(t)\,dt
        \ge
        \frac{\left(\int_0^\infty g(t)\,dt\right)^2}{2\sup g(x)}
\]
with equality if and only if $g$ is constant on an interval $[0,a]$ for some $a \geq 0.$
\end{lemma}

\begin{lemma}[Borell \cite{borell}]
\label{lem:moment-upper-logconcave}
Let $g:[0,\infty)\to[0,\infty)$ be integrable and log-concave. 
Then
\[
        \int_0^\infty t g(t)\,dt
        \le
        \frac{\left(\int_0^\infty g(t)\,dt\right)^2}{g(0)}
\]
with equality if and only if $g$ is exponential on $[0,+\infty).$
\end{lemma}

We first consider the limiting case $\alpha=\infty$.
\begin{theorem} \label{minentropyreverseTh}
Let $f:\mathbb R\to[0,\infty)$ be a log-concave probability density such that
\[
        \int_{\mathbb R} x f(x)\,dx=0.
\]
Then
\[
        h_\infty(f^+)
        \ge
        h_\infty(f)-\log(1+\sqrt2).
\]
Equivalently,
\begin{equation} \label{sqrt2equality}
        \sup_{\mathbb R_+} f(x)
        \le
        (1+\sqrt2)
        \left(\sup_{x\in\mathbb R} f(x)\right)
        \int_0^\infty f(x)\,dx .
\end{equation}
There is equality if and only if $f$ is constant on a finite interval $[-a,0], a \geq 0,$ and exponential on $\mathbb R_+$.
%\textcolor{blue}{this cannot be true, density %constant in $R^{-}$ not integrable} %%
%{\color{red}Lampros: Clearly we mean constant on a %bounded interval}
\end{theorem}

\begin{proof}

Set
\[
        p_{+}:=\int_0^\infty f(x)\,dx
        \quad \text{and } \quad
        p_{-}:=\int_{-\infty}^0 f(x)\,dx.
\]

Note first that if $p_{-}\leq \sqrt{2} p_{+}$ the inequality follows immediately since 
\begin{equation} \label{qleqsqrtp}
    \sup_{\mathbb R_+} f(x)=(p_{+}+p_{-})\sup_{\mathbb R_+} f(x) \leq (1+\sqrt2)p_{+} \sup_{\mathbb R_+} f(x) \leq (1+\sqrt2)p_{+} \sup_{x\in \mathbb R} f(x).
\end{equation}

If $\sup_{\mathbb R_+} f(x) = \sup_{x\in \mathbb R} f(x)$, then we apply Lemma \ref{lem:moment-lower-decreasing} to $g(t) = f(-t)\mathbf 1_{\mathbb R_+}$ and Lemma \ref{lem:moment-upper-logconcave} to $f(t) \mathbf 1_{\mathbb R_+}$ and obtain 
$$
\frac{(p_{-})^{2}}{2f(0)}\leq \frac{(p_{+})^{2}}{f(0)},
$$
whence 
\begin{equation} \label{asabove}
    p_{-} \leq \sqrt{2}p_{+}
\end{equation}
and the claimed inequality holds by \eqref{qleqsqrtp}. 

So it suffices to prove the inequality when $\sup_{x \in \mathbb R_-}f(x) = \sup_{x \in \mathbb R}f(x)$ and $p_{-} > \sqrt2 p_{+}.$ Then by Lemmas \ref{lem:moment-lower-decreasing} and \ref{lem:moment-upper-logconcave} again, 
\begin{equation} \label{lemmasqgreater}
\frac{(p_{-})^{2}}{2\sup_{x \in \mathbb R_-}f(x)}\leq \frac{(p_{+})^{2}}{f(0)}.
\end{equation}
But then, since in this case $\sup_{\mathbb R_+} f(x) = f(0)$
$$
\frac{\sup_{\mathbb R_+} f(x)}{\sup_{x \in \mathbb R}f(x) \int_0^{\infty}f(x)dx}\leq \frac{2p_{+}}{(p_{-})^{2}} = (1+\frac{p_{-}}{p_{+}})\frac{2(p_{+})^{2}}{(p_{-})^{2}} \leq 1+\sqrt{2}
$$
since $\frac{p_{+}}{p_{-}} < \frac{1}{\sqrt{2}}.$ %\textcolor{blue}{typo?} %%{\color{red}Lampros: Yes, fixed it.}

Now suppose there is equality in \eqref{sqrt2equality}. By \eqref{qleqsqrtp}, $p_{-}<\sqrt2 p_{+} $ is impossible as otherwise the inequality would be strict. If $\sup_{\mathbb R_+} f(x) = \sup_{x\in \mathbb R} f(x)$, we must have $p_{-} = \sqrt{2}p_{+}$ as in \eqref{asabove}. On the other hand, if $p_{-} \geq \sqrt 2  p_{+}$ and $\sup_{\mathbb R_-} f(x) = \sup_{x\in \mathbb R} f(x)$ we must have $p_{-} = \sqrt{2}p_{+}$, as otherwise we would have strict inequality by \eqref{lemmasqgreater}. 

So it suffices to characterize the equality case if $p_{-} = \sqrt 2 p_{+}$. But then 
% \eqref{lemmasqgreater} says $$f(0) \leq \sup_{\mathbb R_-} f(x)$$
% and hence $f$ is non-increasing on $\mathbb R_+$
% \textcolor{blue}{smth has to be corrected}
% . Therefore, $$\max_{\mathbb R}f(x) = \max_{\mathbb R_-}f(x),$$
% whereas 
equality in \eqref{sqrt2equality} reads 
$$
\sup_{x\in \mathbb R_+} f(x) = \max_{\mathbb R}f(x).
$$ %Hence there is equality in \eqref{lemmasqgreater}. 
In this case there must be equality in both the applications of Lemma \ref{lem:moment-lower-decreasing} to $g(t)$ and Lemma \ref{lem:moment-upper-logconcave}, which yields the claimed characterization. 
\end{proof}
    
    
    % The barycenter
    % condition gives
    % \[
    % \int_0^\infty x f(x)\,dx
    % =
    % \int_{-\infty}^0 |x|f(x)\,dx .
    % \]
    
    
    
    
    
    % We first record a consequence of Lemmas
    % \ref{lem:moment-lower-decreasing} and \ref{lem:moment-upper-logconcave}.
    
    % \medskip
    
    % \noindent
    % \textbf{Claim.}
    % If
    % \[
    % \sup_{\mathbb R_-} f=f(0),
    % \]
    % then
    % \[
    % q\le \sqrt2\,p.
    % \]
    
    % Indeed, under this assumption the function \(x\mapsto f(-x)\) is
    % non-increasing on \([0,\infty)\). Applying Lemma
    % \ref{lem:moment-lower-decreasing} to \(g(x)=f(-x)\), we get
    % \[
    % \int_{-\infty}^0 |x|f(x)\,dx
    % \ge
    % \frac{q^2}{2f(0)}.
    % \]
    % On the other hand, applying Lemma \ref{lem:moment-upper-logconcave} to
    % \(g(x)=f(x)\), we get
    % \[
    % \int_0^\infty x f(x)\,dx
    % \le
    % \frac{p^2}{f(0)}.
    % \]
    % Using the barycenter condition, we obtain
    % \[
    % \frac{q^2}{2f(0)}
    % \le
    % \int_{-\infty}^0 |x|f(x)\,dx
    % =
    % \int_0^\infty x f(x)\,dx
    % \le
    % \frac{p^2}{f(0)}.
    % \]
    % Therefore
    % \[
    % q\le \sqrt2\,p.
    % \]
    
    % We now prove the theorem by distinguishing three cases.
    
    % \medskip
    
    % \noindent
    % \textbf{Case 1: \(q\le p\).}
    
    % Then
    % \[
    % 1=p+q\le 2p.
    % \]
    % Since
    % \[
    % \sup_{\mathbb R_+} f\le \sup_{\mathbb R} f,
    % \]
    % we have
    % \[
    % \sup_{\mathbb R_+} f
    % \le
    % \sup_{\mathbb R} f
    % \le
    % 2p\sup_{\mathbb R} f
    % \le
    % (1+\sqrt2)p\sup_{\mathbb R} f.
    % \]
    % This is the desired estimate.
    
    % \medskip
    
    % \noindent
    % \textbf{Case 2: \(q>p\) and \(\sup_{\mathbb R_+}f=\sup_{\mathbb R}f\).}
    
    % Since the maximum is attained on the positive side, log-concavity implies
    % \[
    % \sup_{\mathbb R_-} f=f(0).
    % \]
    % By the claim,
    % \[
    % q\le \sqrt2\,p.
    % \]
    % Thus
    % \[
    % 1=p+q\le (1+\sqrt2)p.
    % \]
    % Using \(\sup_{\mathbb R_+}f=\sup_{\mathbb R}f\), we get
    % \[
    % \sup_{\mathbb R_+} f
    % =
    % \sup_{\mathbb R} f
    % \le
    % (1+\sqrt2)p\sup_{\mathbb R} f.
    % \]
    % This proves the desired estimate in the second case.
    
    % \medskip
    
    % \medskip
    
    % \noindent
    % \textbf{Case 3: \(q>p\) and \(\sup_{\mathbb R_+}f<\sup_{\mathbb R}f\).}
    
    % In this case the maximum of \(f\) is attained on the negative side. Hence,
    % by log-concavity, \(f\) is non-increasing on \(\mathbb R_+\), and therefore
    % \[
    % \sup_{\mathbb R_+}f=f(0).
    % \]
    % Set
    % \[
    % c:=f(0),\qquad
    % p:=\int_0^\infty f(x)\,dx,\qquad
    % q:=\int_{-\infty}^0 f(x)\,dx.
    % \]
    % We construct an auxiliary log-concave function \(\widehat f\) of the form
    % \[
    % \widehat f(x):=
    % \begin{cases}
    % 0, & x<-\ell,\\[1mm]
    % c, & -\ell\le x\le a,\\[1mm]
    % c e^{-(x-a)/t}, & x\ge a,
    % \end{cases}
    % \]
    % where \(\ell,a\ge0\) and \(t>0\) are chosen so that
    % \[
    % \int_{\mathbb R_+}\widehat f=\int_{\mathbb R_+}f=p,
    % \qquad
    % \int_{\mathbb R}\widehat f=\int_{\mathbb R}f=1,
    % \qquad
    % \int_{\mathbb R}x\widehat f(x)\,dx=0.
    % \]
    
    % Let us compute the parameters. First,
    % \[
    % \int_{\mathbb R_+}\widehat f
    % =
    % c a+\int_a^\infty c e^{-(x-a)/t}\,dx
    % =
    % c(a+t).
    % \]
    % Thus we impose
    % \[
    % a+t=\frac pc.
    % \]
    % Moreover,
    % \[
    % \int_{\mathbb R}\widehat f
    % =
    % c(\ell+a+t),
    % \]
    % and since \(\int_{\mathbb R}\widehat f=1=p+q\), we get
    % \[
    % \ell+a+t=\frac{p+q}{c}.
    % \]
    % Combining this with \(a+t=p/c\), we obtain
    % \[
    % \ell=\frac qc.
    % \]
    
    % It remains to impose the barycenter condition. We have
    % \[
    % \int_{\mathbb R}x\widehat f(x)\,dx
    % =
    % c\int_{-\ell}^a x\,dx
    % +
    % c\int_a^\infty x e^{-(x-a)/t}\,dx.
    % \]
    % Hence
    % \[
    % \int_{\mathbb R}x\widehat f(x)\,dx
    % =
    % c\left(
    % \frac{a^2-\ell^2}{2}
    % +
    % at+t^2
    % \right).
    % \]
    % Therefore the condition
    % \[
    % \int_{\mathbb R}x\widehat f(x)\,dx=0
    % \]
    % is equivalent to
    % \[
    % \frac{a^2-\ell^2}{2}+at+t^2=0.
    % \]
    
    % Since \(a+t=p/c\) and \(\ell=q/c\), writing
    % \[
    % B:=\frac pc,
    % \qquad
    % L:=\frac qc,
    % \]
    % we have
    % \[
    % a+t=B,
    % \qquad
    % \ell=L.
    % \]
    % Thus \(t=B-a\), and the barycenter equation becomes
    % \[
    % a^2-2Ba+2B^2-L^2=0.
    % \]
    % Equivalently,
    % \[
    % a
    % =
    % B-\sqrt{L^2-B^2},
    % \qquad
    % t
    % =
    % \sqrt{L^2-B^2}.
    % \]
    % Therefore the construction is admissible whenever
    % \[
    % B\le L\le \sqrt2\,B,
    % \]
    % that is,
    % \[
    % p\le q\le \sqrt2\,p.
    % \]
    
    % With these choices, \(\widehat f\) is log-concave and satisfies
    % \[
    % \int_{\mathbb R_+}\widehat f=\int_{\mathbb R_+}f,
    % \qquad
    % \int_{\mathbb R}\widehat f=\int_{\mathbb R}f,
    % \qquad
    % \int_{\mathbb R}x\widehat f(x)\,dx=0.
    % \]
    % Moreover,
    % \[
    % \sup_{\mathbb R_+}\widehat f
    % =
    % \sup_{\mathbb R}\widehat f
    % =
    % c.
    % \]
    % Thus \(\widehat f\) falls under the case already proved, namely the case in
    % which the maximum is attained on the positive side. Applying that case to
    % \(\widehat f\), we obtain
    % \[
    % \int_{\mathbb R}\widehat f
    % \le
    % (1+\sqrt2)\int_{\mathbb R_+}\widehat f.
    % \]
    % Using the identities defining \(\widehat f\), this gives
    % \[
    % \int_{\mathbb R} f
    % \le
    % (1+\sqrt2)\int_{\mathbb R_+} f.
    % \]
    % Since \(f\) is a probability density,
    % \[
    % 1\le (1+\sqrt2)p.
    % \]
    % Therefore
    % \[
    % \sup_{\mathbb R_+} f
    % \le
    % \sup_{\mathbb R} f
    % \le
    % (1+\sqrt2)p\,\sup_{\mathbb R} f.
    % \]
    % This is the desired estimate in this subcase.
    
    % \medskip
    
    % \noindent
    % \textbf{Case 4: \(q>\sqrt2\,p\) and
    % \(\sup_{\mathbb R_+}f<\sup_{\mathbb R}f\).}
    
    % For convenience, write
    % \[
    % c:=f(0)=\sup_{\mathbb R_+}f,
    % \qquad
    % M:=\sup_{\mathbb R}f .
    % \]
    % Since the maximum of \(f\) is attained on the negative side, log-concavity
    % implies that \(f\) is non-increasing on \(\mathbb R_+\), and hence indeed
    % \[
    % \sup_{\mathbb R_+}f=f(0)=c.
    % \]
    % By Lemma \ref{lem:moment-upper-logconcave} 
    % we obtain
    % \[
    % \int_0^\infty x f(x)\,dx
    % \le
    % \frac{p^2}{c},
    % \]
    % and by Lemma \ref{lem:moment-lower-decreasing} \[
    % \int_{-\infty}^0 (-x)f(x)\,dx
    % \ge
    % \frac{q^2}{2M}.
    % \]
    
    % Combining the two moment estimates with the barycenter condition gives
    % \[
    % \frac{q^2}{2M}
    % \le
    % \int_{-\infty}^0 (-x)f(x)\,dx
    % =
    % \int_0^\infty x f(x)\,dx
    % \le
    % \frac{p^2}{c}.
    % \]
    % Hence
    % \[
    % \frac{c}{M}
    % \le
    % \frac{2p^2}{q^2}\leq \frac{(1+\sqrt{2})p}{p+q},
    % \]
    % since $q>\sqrt{2}p$. This finishes the proof.
    
    % \end{proof}

% \begin{comment}
    
% \begin{lemma}[The subcase \(q>\sqrt2\,p\)]
% Let \(f\) be a log-concave probability density on \(\mathbb R\) with barycenter
% zero. Set
% \[
%         p:=\int_0^\infty f(x)\,dx,
%         \qquad
%         q:=\int_{-\infty}^0 f(x)\,dx,
% \]
% so that \(p+q=1\). Assume that
% \[
%         q>\sqrt2\,p
% \]
% and that
% \[
%         \sup_{\mathbb R_+} f < \sup_{\mathbb R} f.
% \]
% Then
% \[
%         \sup_{\mathbb R_+} f
%         \le
%         (1+\sqrt2)
%         \left(\int_0^\infty f(x)\,dx\right)
%         \sup_{\mathbb R} f .
% \]
% \end{lemma}

% \begin{proof}
% Write
% \[
%         c:=f(0)=\sup_{\mathbb R_+}f,
%         \qquad
%         M:=\sup_{\mathbb R}f .
% \]
% Since the maximum of \(f\) is attained on the negative side, log-concavity
% implies that \(f\) is non-increasing on \(\mathbb R_+\), and hence indeed
% \[
%         \sup_{\mathbb R_+}f=f(0)=c.
% \]

% Since \(f\) has barycenter zero,
% \[
%         \int_0^\infty x f(x)\,dx
%         =
%         \int_{-\infty}^0 (-x) f(x)\,dx .
% \]

% We first estimate the positive first moment from above. Define
% \[
%         S(x):=\int_x^\infty f(u)\,du,
%         \qquad x\ge0.
% \]
% The function \(S\) is log-concave on \(\mathbb R_+\). Moreover,
% \[
%         S(0)=p
%         \qquad \text{ and } \quad 
%         S'(0)=-f(0)=-c .
% \]
% Therefore, by concavity of \(\log S\),
% \[
%         S(x)
%         \le
%         S(0)\exp\left(\frac{S'(0)}{S(0)}x\right)
%         =
%         p\exp\left(-\frac{c}{p}x\right).
% \]
% Since,
% \[
%         \int_0^\infty x f(x)\,dx
%         =
%         \int_0^\infty S(x)\,dx,
% \]
% we obtain
% \[
%         \int_0^\infty x f(x)\,dx
%         \le
%         \int_0^\infty p\exp\left(-\frac{c}{p}x\right)\,dx
%         =
%         \frac{p^2}{c}.
% \]

% We now estimate the negative first moment from below. Since \(0\le f\le M\),
% we claim that
% \[
%         \int_{-\infty}^0 (-x)f(x)\,dx
%         \ge
%         \frac{q^2}{2M}.
% \]
% Indeed, put
% \[
%         h(y):=f(-y),
%         \qquad y\ge0.
% \]
% Then \(0\le h\le M\) and
% \[
%         \int_0^\infty h(y)\,dy=q.
% \]
% Let
% \[
%         A:=\frac qM.
% \]
% Then
% \[
% \begin{aligned}
%         \int_0^\infty y h(y)\,dy-\frac{q^2}{2M}
%         &=
%         \int_0^\infty y h(y)\,dy
%         -
%         M\int_0^A y\,dy                                                     \\
%         &=
%         \int_0^A y(h(y)-M)\,dy
%         +
%         \int_A^\infty y h(y)\,dy .
% \end{aligned}
% \]
% On \([0,A]\), we have \(h(y)-M\le0\) and \(y\le A\), so
% \[
%         y(h(y)-M)\ge A(h(y)-M).
% \]
% On \([A,\infty)\), we have \(y\ge A\), so
% \[
%         yh(y)\ge Ah(y).
% \]
% Therefore
% \[
% \begin{aligned}
%         \int_0^\infty y h(y)\,dy-\frac{q^2}{2M}
%         &\ge
%         A\int_0^A (h(y)-M)\,dy
%         +
%         A\int_A^\infty h(y)\,dy                                      \\
%         &=
%         A\left(\int_0^\infty h(y)\,dy-MA\right)                       \\
%         &=
%         A(q-MA)
%         =
%         0.
% \end{aligned}
% \]
% Thus
% \[
%         \int_{-\infty}^0 (-x)f(x)\,dx
%         =
%         \int_0^\infty y h(y)\,dy
%         \ge
%         \frac{q^2}{2M}.
% \]

% Combining the two moment estimates with the barycenter condition gives
% \[
%         \frac{q^2}{2M}
%         \le
%         \int_{-\infty}^0 (-x)f(x)\,dx
%         =
%         \int_0^\infty x f(x)\,dx
%         \le
%         \frac{p^2}{c}.
% \]
% Hence
% \[
%         \frac{c}{M}
%         \le
%         \frac{2p^2}{q^2}.
% \]
% Since \(p+q=1\), this becomes
% \[
%         \frac{c}{M}
%         \le
%         \frac{2p^2}{(1-p)^2}.
% \]
% It remains to compare the right-hand side with \((1+\sqrt2)p\). The
% assumption \(q>\sqrt2\,p\) is equivalent to
% \[
%         1-p>\sqrt2\,p,
% \]
% and hence
% \[
%         p<\frac{1}{1+\sqrt2}.
% \]
% For such \(p\),
% \[
%         \frac{2p^2}{(1-p)^2}
%         \le
%         (1+\sqrt2)p.
% \]
% Indeed, this inequality is equivalent to
% \[
%         2p\le (1+\sqrt2)(1-p)^2,
% \]
% or 
% $$
% (1+\sqrt{2})\Bigl(p-\frac{1}{1+\sqrt2}\Bigr)\Bigl(p-(1+\sqrt2)\Bigr) \geq 0,
% $$
% which holds for \(0\le p\le 1/(1+\sqrt2)\), with equality at
% \(p=1/(1+\sqrt2)\).

% Therefore
% \[
%         \frac{c}{M}
%         \le
%         (1+\sqrt2)p.
% \]
% Recalling that \(c=\sup_{\mathbb R_+}f\), \(M=\sup_{\mathbb R}f\), and
% \(p=\int_0^\infty f\), we obtain
% \[
%         \sup_{\mathbb R_+} f
%         \le
%         (1+\sqrt2)
%         \left(\int_0^\infty f(x)\,dx\right)
%         \sup_{\mathbb R} f .
% \]
% This proves the claim.
% \end{proof}
% \end{comment}



% To verify it, note that
% \[
% h(X)
% =
% -\int_{\mathbb{R}^n} f(x)\log f(x)\,dx
% =
% -\int_{H^+} f(x)\log f(x)\,dx
% -
% \int_{H^-} f(x)\log f(x)\,dx.
% \]
% On $H^+$ we have $f(x)=p^+ f^+(x)$, hence
% \[
% -\int_{H^+} f(x)\log f(x)\,dx
% =
% -\int_{H^+} p^+ f^+(x)\log\big(p^+ f^+(x)\big)\,dx
% =
% p^+ h(X^+) - p^+\log p^+.
% \]
% Similarly,
% \[
% -\int_{H^-} f(x)\log f(x)\,dx
% =
% p^- h(X^-) - p^-\log p^-.
% \]

Next we prove the $\alpha = 1$ case of Theorem \ref{secondThIntro}. In the proof we use Proposition \ref{prop:normalized-variational-inequality}, which is proved at the end of this section. 
\begin{theorem}
%[Sharp Gr\"unbaum--type inequality for differential entropy]
\label{thm:sharp-reverse-grunbaum-shannon}
Let \(X\) be a real-valued random variable with log-concave density \(g\), and assume that
\[
        \int_{\mathbb R} xg(x)\,dx=0.
\]
Then
\[
        h(g^+)
        \ge
        h(g)-\frac{e}{e-1}H_2(1/e).
\]
Equality holds if and only if $g$ is the exponential density
\[
        g(x)=\lambda e^{\lambda(x-1/\lambda)}
        \mathbf 1_{(-\infty,1/\lambda]}(x),
        \qquad \lambda>0.
\]
\end{theorem}

%\textcolor{blue}{It s not $p+m+=p-m-$}

\begin{proof}
 
Let
\[
        p_+ := \int_0^\infty g(x)\,dx,
        \quad p_- = 1-p_+ \quad \text{and }
        g^+(x):=\frac{g(x)\mathbf 1_{\{x\ge0\}}}{p_+}, g^-(x):=\frac{g(x)\mathbf 1_{\{x<0\}}}{p_-}
\]
be the densities of \(X^+\) and \(X^-\). Let also
\[
        m_+ := \int_0^\infty x g^+(x)\,dx \quad \text{and } m_- = \int_{-\infty}^0 |x| g^-(x)\,dx.
\]
Therefore $p_+m_+ = p_-m_-$. We have
$$
h(g) = H_2(p_+) + p_+h(g^+)+p_-h(g^-)
$$
and thus 
\begin{align}
    h(g)-h(g^+) &= H_2(p_+) + p_-\bigl(h(g^-) - h(g^+)\bigr).
\end{align}
Since the negative exponential maximizes entropy under fixed mean on $\mathbb R_-$ \cite[Example 12.2.5]{cover:book} (as can be seen by $D(g^-\|\mathrm{exp}_-(\frac{1}{m_-})) \geq 0$), we obtain after applying the second part of Proposition \ref{prop:normalized-variational-inequality} to $g^+$,
\begin{align}
    h(g)-h(g^+) &\leq H_2(p_+) + p_-\bigl( 1 + \log{{m_-}} - \log{m_+}-2\log{(e-1)} +\frac{1}{e-1}\bigr) \\
    &= H_2(p_+) + p_-\bigl( 1 -\log{p_-} + \log{p_+}-2\log{(e-1)} +\frac{1}{e-1}\bigr) \\ \label{maximizingexpr}
    &= 2H_2(p_-) + p_-\bigl( 1 -2\log{(e-1)} +\frac{1}{e-1}\bigr) + \log{(1-p_-)}.
\end{align}
Considering the function $x \mapsto -2x\log{x} -2(1-x)\log{(1-x)} + cx + \log{(1-x)}, x \in [0,1],$ where $c = 1-2\log{(e-1)} +\frac{1}{e-1},$ and setting its derivative to zero, we see that the maximum is achieved at $x = 1/e$. Therefore, \eqref{maximizingexpr} is maximized for $p_- = 1/e$ and therefore
\begin{align}
h(g)-h(g^+) &\leq 2H_2(1/e) + 1/e\bigl( 1 -2\log{(e-1)} +\frac{1}{e-1}\bigr) + \log{(1-1/e)} \\
&= \frac{e}{e-1}H_2(1/e)
\end{align}
after a short calculation, with equality if and only if $p_- = 1/e$. The equality case follows from the equality case in the functional form of Gr{\"u}nbaum's inequality for log-concave densities, see for example~\cite[Theorem 1.4.]{fradelizi2026grunbaum}.
%  Applying that proposition gives
% \[
%         h(g^+)
%         =
%         h(\bar f)+\log m_+
%         \ge
%         \log m_+
%         +
%         2\log(e-1)-\frac1{e-1}.
% \]
% Since \(m_+=A/p_+\), this becomes
% \[
%         h(g^+)
%         \ge
%         \log A-\log p_+
%         +
%         2\log(e-1)-\frac1{e-1}.
% \]
% On the other hand,  the centered density \(g\) satisfies (\textcolor{red}{HERE DOUBT-!}),
% \[
%         h(g)\le 2+\log A.
% \]
% Combining the last two inequalities and using Gr\"unbaum's inequality
% \[
%         p_+\le \frac{e-1}{e},
% \]
% we get
% \[
%         h(g)-h(g^+)
%         \le
%         1-\log(e-1)+\frac1{e-1}
%         =
%         \frac{e}{e-1}H_2(1/e).
% \]
% Therefore
% \[
%         h(g^+)
%         \ge
%         h(g)-\frac{e}{e-1}H_2(1/e),
% \]
% which proves the inequality in Theorem
% \ref{thm:sharp-reverse-grunbaum-shannon}, assuming Proposition
% \ref{prop:normalized-variational-inequality}.
\end{proof}

Next we establish the lower bound on the entropy of a density on the positive reals that was used in the proof of Theorem \ref{thm:sharp-reverse-grunbaum-shannon}, but which may also be of independent interest. It implies in particular that the entropy of the restriction of a centered log-concave density on $\mathbb R_+$ is minimized by an exponential; this is analogous to \cite{MNR} where the entropy among log-concave densities with fixed variance is minimized for an exponential.
\begin{proposition}
\label{prop:normalized-variational-inequality}
Let \(f\) be a log-concave probability density on \(\mathbb R_+\) such that $f(0)>0$ and let $\kappa:=(\log f)'_+(0)$.

If 
$$
        \sqrt{f(0)}\ge \kappa\sqrt{\int_0^\infty tf(t) \ dt},
$$
then
\begin{equation} \label{variationalineq}
        h(f)\ge \log\Biggl({\int_0^\infty tf(t)\dd t}\Biggr)+2\log(e-1)-\frac1{e-1},
\end{equation}
where the right hand side is the entropy of the truncated exponential $f^*(x) = \frac1{\lambda(e-1)^2}e^{\frac{x}{\lambda (e-1)}}\1_{(0,\lambda (e-1))}(x), $ with $\lambda = \int_0^\infty tf(t)\dd t.$

In particular, if $g$ is a centered log-concave density on \(\mathbb R\), i.e., 
$
\int_{\mathbb R} t g(t)\dd t=0,
$
then \eqref{variationalineq} holds with $g^+$ in place of $f$, where
$
g^+(t):=\frac{g(t)}{\int_0^\infty g(x)\dd x}\1_{\{t>0\}}.
$ 
\end{proposition}

The proof is divided in several lemmas.

\begin{lemma} \label{lemmasatisfiescond}
Let $f$ be a centered log-concave density. If
$$
\frac{d}{dx} \log{f(x)}\Bigg|_{0^+}= \bigl(\log{f}\bigr)_+'(0) \geq 0,
$$
then 

\[
f(0)\ge {\Bigl(\bigl(\log{f}\bigr)_+'(0)\Bigr)^2}{\int_0^\infty tf(t) \ dt}.
\]
\end{lemma}

\begin{proof}


Since \(f\) is log-concave, for \(t<0\),

\[
\log f(t)
\le
\log f(0)+\bigl(\log{f}\bigr)_+'(0)t.
\]


Because \(f\) is centered,

\[
\int_{-\infty}^0 (-t){f(t)}\dd t
=
\int_0^\infty t{f(t)}\dd t
.
\]
Since \(\bigl(\log{f}\bigr)_+'(0) \geq 0\),

\[
\int_0^\infty t{f(t)}\dd t \leq f(0)\int_{-\infty}^0 (-t)e^{\bigl(\log{f}\bigr)_+'(0)t}\dd t
=
\frac{f(0)}{{\bigl(\bigl(\log{f}\bigr)_+'(0)\bigr)^{2}}
}\]
and the claim follows.
\end{proof}

\begin{comment}
\begin{lemma} \label{lem:existence-compact-problem}
    Let $R>1$ and $\mathcal{P}(R)$ be the set of measures on $[0,R]$ that have log-concave densities which are at most three-piece log-affine, have mean $1$, satisfy $f(0)>0$ and
\begin{equation}\label{condition}
        \sqrt{f(0)}\ge (\log f)'_+(0).
\end{equation}
 Then $\mathcal{P}(R)$ is compact and entropy is a continuous functional on $\mathcal{P}(R)$  with respect to weak convergence. 

\end{lemma}
\begin{proof} Let $(\mu_n)_n$ be a sequence of probability measures in $\mathcal{P}(R)$ which converges weakly to a measure $\mu$. Denote by $f_n$  the density of $\mu_n$. 
First, we observe that the set of log-concave probability measures on a compact set is compact with respect to weak convergence, hence $\mu$ is log-concave on $[0,R]$. Thus $\mu$ is either a Dirac measure or has a log-concave density. If we write $f_n(x) = e^{-V_n(x)}$, then by convexity, condition \eqref{condition} is equivalent to 
  \begin{equation} \label{vequiv}
      V_n(x) \geq V_n(0) - e^{-V_n(0)/2}x \quad \text{ for every } x \in [0,R].
  \end{equation}
 
  
  Condition \eqref{vequiv} implies that $\mu$ cannot be Dirac (since otherwise $V_n(1) \to -\infty$) and therefore must have a continuous log-concave density $f$. Moreover, $f_n \to f$ a.e. Also, $f = e^{-V}$ also satisfies \eqref{vequiv}, has mean $1$ and is three-piece log-affine. Hence $\mathcal{P}(R)$ is closed and it is also compact by tightness of measures.

  Now $f_n \to f$ a.e.  and without loss of generality we assume that $f_n$ and $f$ are continuous on their supports, say $S_n$ and $S$ respectively. Since $f_n$ is log-concave it is classical that $f_n \to f$ uniformly on $S$ (see for example %and $f_n \to 0$ uniformly on compact subsets of $S^C$
  \cite[Fact 2.5]{corderofradelizi}). Therefore $f_n(0) \to f(0)$. Condition \eqref{vequiv} is equivalent to 
  $$
  f_n(x) \leq f_n(0)e^{\sqrt{f_n(0)}x}, \qquad x \in [0,R]
  $$
  and thus $f_n(0)$ is bounded. By dominated convergence 
  $$
  h(f_n) \to h(f). 
  $$
  % Since on $[0,R]$, $h(f) = \log{R}-D(f\|\mathrm{Uni}([0,R]))$ and relative entropy is lower-semicontinuous with respect to weak convergence, entropy is upper-semicontinuous. 
  % So it suffices to show that if $f_n \to f$ a.e. then 
  % \begin{equation} \label{liminfsuff}
  %     \liminf_nh(f_n) \geq h(f).
  % \end{equation}
  % But 
  % $$
  % h(f_n) = \int_{\{f_n \leq 1\}}-f_n \log{f_n} - \int_{\{f_n > 1\}}f_n \log{f_n}.
  % $$
  % The first integral satisfies
  % \begin{equation} \label{firstint}
  % \liminf_n \int_{\{f_n \leq 1\}}-f_n \log{f_n} \geq \int_{\{f \leq 1\}}-f \log{f}
  % \end{equation}
  % by Fatou. 
  % For the second integral we note that the sets $\{f_n > 1\}$ are convex for every $n$ by log-concavity and since they have finite Lebesgue measure they are uniformly bounded. Therefore by bounded convergence \textcolor{blue}{here we need $\sup_n \mid \mid f_n\mid\mid_{infty}<\infty$ no? }
  %  \begin{equation} \label{secondint}
  % \int_{\{f_n > 1\}}f_n \log{f_n} \to \int_{\{f > 1\}}f \log{f}.
  % \end{equation}
  % The result follows since \eqref{firstint} and \eqref{secondint} give \eqref{liminfsuff}.
\end{proof}
% \begin{lemma}[Existence in the compact problem]
% \label{lem:existence-compact-problem}
% For every \(R>1\), the entropy functional attains its minimum on \(\mathcal F_R\).
% \end{lemma}

% \begin{proof}
% Let \((f_n)\subset \mathcal F_R\) be a minimizing sequence. Since all the measures
% \(f_n(x)\,dx\) are supported on the compact interval \([0,R]\), we may extract a
% subsequence converging weakly to a probability measure \(\mu\).

% The mass and mean constraints pass to the limit, hence
% \[
%         \int_0^R d\mu=1,
%         \qquad
%         \int_0^R x\,d\mu(x)=1.
% \]
% We claim that \(\mu\) is non-degenerate. Indeed, if the sequence converged to a Dirac
% mass, the mean constraint would force the limit to be \(\delta_1\). Such a
% concentration is incompatible with the endpoint condition: for a log-concave density
% concentrating at \(1\), the value at the origin tends to \(0\), while a positive initial
% logarithmic slope must become unbounded. This contradicts
% \[
%         f_n(0)\ge \big((\log f_n)'_+(0)\big)^2
% \]
% whenever \((\log f_n)'_+(0)>0\). \textcolor{red}{perhaps there is a little gap here if it the function is zero until almost 1... }

% Thus \(\mu\) is a non-degenerate log-concave probability measure on \([0,R]\).
% Therefore it has a log-concave density \(f\). By the standard compactness properties
% of one-dimensional log-concave densities, after passing to a further subsequence,
% \[
%         f_n\to f
%         \qquad \text{a.e. on }[0,R],
% \]
% and
% \[
%         \sup_n\|f_n\|_\infty<\infty.
% \]
% Moreover, the endpoint condition is stable under this non-degenerate log-concave
% convergence, and hence \(f\in\mathcal F_R\).

% It remains to pass to the limit in the entropy. Let
% \[
%         \varphi(t):=t\log\frac1t.
% \]
% Since the sequence \((f_n)\) is uniformly bounded in \(L^\infty\), the functions
% \(\varphi(f_n)\) are dominated by an integrable constant on \([0,R]\). Hence, by
% dominated convergence,
% \[
%         h(f_n)
%         =
%         \int_0^R \varphi(f_n(x))\,dx
%         \longrightarrow
%         \int_0^R \varphi(f(x))\,dx
%         =
%         h(f).
% \]
% Since \((f_n)\) was minimizing, we obtain
% \[
%         h(f)=\inf_{\mathcal F_R}h.
% \]
% Therefore the minimum is attained.
% \end{proof}

We next record the first structural reduction. It is a standard degree-of-freedom
argument in the sense of \cite[Section~2]{fradeliziguedon};
in dimension one, the relevant characterization is
\cite[Proposition~2]{fradeliziguedon}.

\begin{proposition}[Reduction to three pieces]
\label{prop:reduction-three-pieces}
For \(R>1\), let
\(\mathcal F_R\) be the class of log-concave probability densities \(f=e^u\)
supported on \([0,R]\) such that
\[
        \int_0^R f(x)\,dx=1,
        \qquad
        \int_0^R x f(x)\,dx=1,
\]
and such that, whenever
\[
        \kappa:=u'_+(0)>0,
\]
one has
\[
        f(0)\ge \kappa^2.
\]
Let \(f\in\mathcal F_R\) be a minimizer of \(h\) over \(\mathcal F_R\). Then
\(\log f\) is affine on at most three intervals.
\end{proposition}

\begin{proof}
Write \(f=e^{-V}\), with \(V\) convex. Assume, by contradiction, that \(f\) is not
three-piece log-affine. By the one-dimensional degree-of-freedom construction of \cite[Proof of Proposition~2]{fradeliziguedon}, there are enough
independent admissible perturbation directions, all vanishing on an initial interval.
Taking a non-trivial linear combination, we may choose a bounded continuous function
\(W\), still vanishing near \(0\), such that, for all sufficiently small \(t>0\), 
\[
        f_\pm:=f(1\pm tW)
\]
are log-concave and
\[
        \int_0^R Wf=0,
        \qquad
        \int_0^R xWf=0.
\]
The two identities preserve mass and mean, while the fact that \(W\) vanishes near
\(0\) preserves the endpoint condition. Hence
\[
        f_+,f_-\in\mathcal F_R.
\]

Finally,
\[
        f=\frac{f_++f_-}{2}.
\]
Since \(x\mapsto x\log x\) is strictly convex and \(W\neq0\),
\[
        h(f)>
        \frac12 h(f_+)+\frac12 h(f_-).
\]
Thus at least one of \(f_+\) and \(f_-\) has entropy strictly smaller than \(h(f)\),
contradicting the minimality of \(f\). Therefore \(\log f\) is affine on at most three
intervals.
\end{proof}
\end{comment}



\begin{proposition}%[Existence and reduction to three pieces]
\label{prop:existence-reduction}
Let \(R>1\), and let \(\mathcal F_R\) be the class of log-concave
probability densities \(f=e^u\) supported on \([0,R]\) such that
\[
\int_0^R f(x)\,dx=1,
\qquad
\int_0^R xf(x)\,dx=1,
\]
and such that
\[
f(0)>0
\qquad\text{and}\qquad
\sqrt{f(0)}\geq u'_+(0).
\]
Then \(\mathcal F_R\) is compact with respect to weak convergence of the associated probability measures,
and entropy is a continuous functional on \(\mathcal F_R\). Consequently, the entropy functional attains its minimum on
\(\mathcal F_R\). Moreover, if \(f\in\mathcal F_R\) is a minimizer,
then \(\log f\) is affine on at most three intervals.
\end{proposition}

\begin{proof}


Let $(\mu_n)_n$ be a sequence of probability measures in $\mathcal F_R $  which converges weakly to a measure $\mu$. Denote by $f_n$  the density of $\mu_n$. 
First, we observe that the set of log-concave probability measures on a compact set is compact with respect to weak convergence, hence $\mu$ is log-concave on $[0,R]$. Thus $\mu$ is either a Dirac measure or has a log-concave density.
The boundary condition at the origin is equivalent to
\[
f_n(x)\leq f_n(0)e^{\sqrt{f_n(0)}\,x},
\qquad x\in[0,R].
\]
The normalization $\int_0^R f_n=1$ prevents $f_n(0)$ from tending to
zero. Moreover, the condition $\int_0^R xf_n(x)\,dx=1$ forces a
uniformly positive amount of mass to lie away from the origin, and
log-concavity then prevents $f_n(0)$ from tending to infinity.
Hence $(f_n)_n$, is uniformly bounded. It follows that the weak limit $\mu$ cannot be a Dirac mass, and
therefore has a log-concave density $f$ with mean $1$. After passing
to a subsequence, assume that $f_n(0)\to l>0$. By
\cite[Fact~2.5]{corderofradelizi}, $f_n\to f$ uniformly on compact
subsets of $\operatorname{int}(\operatorname{supp}f)$. Passing to the
limit in the upper bound above and in the corresponding lower bound
given by log-concavity, and then letting $x\downarrow0$, yields
$l=f(0)$. Hence $f$ also satisfies the boundary condition at the origin, so
$f\in\mathcal F_R$. Thus $\mathcal F_R$ is weakly closed, hence
compact. Finally, the uniform bound on $(f_n)_n$ and dominated
convergence give
\[
h(f_n)\longrightarrow h(f).
\]


\begin{comment}

If we write $f_n(x) = e^{-V_n(x)}$, then by convexity, the endpoint condition  is equivalent to 
  \begin{equation} \label{vequiv}
      V_n(x) \geq V_n(0) - e^{-V_n(0)/2}x \quad \text{ for every } x \in [0,R].
  \end{equation}

 
  
  Condition \eqref{vequiv} implies that $\mu$ cannot be Dirac (since otherwise $V_n(1) \to -\infty$) and therefore must have a continuous log-concave density $f$. Moreover, $f_n \to f$ a.e. Also, $f = e^{-V}$ also satisfies \eqref{vequiv}, has mean $1$. Hence $\mathcal F_R$  is closed and it is also compact by tightness of measures.

  Now $f_n \to f$ a.e.  and without loss of generality we assume that $f_n$ and $f$ are continuous on their supports, say $S_n$ and $S$ respectively. Since $f_n$ is log-concave it is classical that $f_n \to f$ uniformly on $S$ (see for example %and $f_n \to 0$ uniformly on compact subsets of $S^C$
  \cite[Fact 2.5]{corderofradelizi}). 
  \textcolor{blue}{Precisely, on the fact 2.5 they say uniform convergence on any compact subset on the interior of the support }
  Therefore $f_n(0) \to f(0)$. Condition \eqref{vequiv} is equivalent to 
  $$
  f_n(x) \leq f_n(0)e^{\sqrt{f_n(0)}x}, \qquad x \in [0,R]
  $$
  and $f_n(0)$ is a bounded sequence since it is convergent. By dominated convergence 
  $$
  h(f_n) \to h(f). 
  $$
 \end{comment}

We now turn to the structural characterization of the minimizers.
Let \(f\in\mathcal F_R\) be a minimizer and write \(f=e^{-V}\), with
\(V\) convex. 
Assume, by contradiction, that \(f\) is not
three-piece log-affine. By the one-dimensional degree-of-freedom
construction of
\cite[Proof of Proposition~2]{fradeliziguedon}, there are enough
independent admissible perturbation directions, all vanishing on an
initial interval. Taking a non-trivial linear combination, we may
choose a bounded continuous function \(W\), still vanishing near
\(0\), such that, for all sufficiently small \(t>0\),
\[
f_\pm:=f(1\pm tW)
\]
are log-concave and
\[
\int_0^R Wf=0,
\qquad
\int_0^R xWf=0.
\]
The two identities preserve mass and mean, while the fact that \(W\)
vanishes near \(0\) preserves the endpoint condition. Hence
\[
f_+,f_-\in\mathcal F_R.
\]

Finally,
\[
f=\frac{f_++f_-}{2}.
\]
Since \(x\mapsto x\log x\) is strictly convex and \(W\neq0\),
\[
h(f)>
\frac12h(f_+)+\frac12h(f_-).
\]
Thus at least one of \(f_+\) and \(f_-\) has entropy strictly smaller
than \(h(f)\), contradicting the minimality of \(f\). Therefore
\(\log f\) is affine on at most three intervals.
\end{proof}


We will use the following optimization lemma to show that in fact a minimizer must be one-piece log-affine. 
\begin{lemma}[Fritz--John necessary conditions; Proposition 3.3.5 in \cite{bertsekas}] \label{FJ}
Let 

\[
F,G_1,\dots,G_r,C_1,\dots,C_s\in C^1(\mathbb R^n).
\]
Suppose \(z_0\in \mathbb{R}^n\) is a local minimizer of \(F\) subject to

\[
G_i(z)=0,\qquad i=1,\dots,r,
\]
and

\[
C_j(z)\le0,\qquad j=1,\dots,s.
\]
Then there exist multipliers

\[
\tau\ge0,\qquad
\alpha_1,\dots,\alpha_r\in\mathbb R,
\qquad
\gamma_1,\dots,\gamma_s\ge0,
\]
not all zero, such that

\[
\tau\nabla F(z_0)
+
\sum_{i=1}^r \alpha_i\nabla G_i(z_0)
+
\sum_{j=1}^s \gamma_j\nabla C_j(z_0)=0.
\]
Moreover,

\[
\gamma_j C_j(z_0)=0,
\qquad j=1,\dots,s.
\]
\end{lemma}

We are finally ready to give the proof of Proposition \ref{prop:normalized-variational-inequality}.
\begin{proof} [Proof of Proposition \ref{prop:normalized-variational-inequality}.]
    
First note that, by approximating $f$ with compactly supported densities, we may assume that $f$ has a compact support. Then, we observe that, by considering the mean-normalized density  
$ 
\mu f(\mu x), \quad x>0,
$
where $\mu = {\int_0^{\infty}f(x)xdx},$
and using the scaling property of differential entropy, it suffices to prove the normalized inequality
\[
h(f)\ge 2\log(e-1)-\frac1{e-1}
\]
for any log-concave density $f$ on $[0,R]$ with $\int_{\mathbb R_+}f(x)xdx = 1$ and $\sqrt{f(0)} \geq \kappa$. 
By Proposition \ref{prop:existence-reduction}, the entropy functional attains its minimum on
\(\mathcal F_R\) at some density \(f_R\), and every such minimizer is
three-piece log-affine. Therefore, it suffices to prove the normalized
inequality for a three-piece log-affine minimizer. Renaming this
minimizer \(f\), it can be parametrized as
\[
f(x)=e^{u(x)}\1_{(0,R)}(x),
\]
where \(R\in(0,\infty)\) and
\[
u(x)=a+\kappa x-\lambda_1(x-r_1)_+-\lambda_2(x-r_2)_+.
\]
Here
$
0<r_1<r_2<R$ and 
$\lambda_1,\lambda_2\ge0.
$
The slopes of \(u\) are $\kappa,
\kappa-\lambda_1$ and $\kappa-\lambda_1-\lambda_2.$
In particular, concavity is exactly the condition
\[
\lambda_1,\lambda_2\ge0.
\]
There is a change of slope at a point \(r_j\) precisely when
\[
\lambda_j>0.
\]
We will use the Fritz-John necessary conditions to show that if a density of this form is a local minimizer of the entropy functional, then it must be the case that $\lambda_j = 0$ for $j=1,2$ and thus any local minimizer is one-piece log-affine. 

We are minimizing locally the entropy 
\[
h(f)=-\int_0^R u(x)e^{u(x)}\dd x.
\]
in the variables $(a, \kappa, \lambda_1,\lambda_2,r_1, r_2) \in \mathbb R^6$ subject to the constraints
\[
\int_0^R e^{u(x)}\dd x=1,
\quad 
\int_0^R x e^{u(x)}\dd x=1 \quad \text{and } \quad \kappa-e^{\frac{a}{2}}\le0.
\]
We will assume for contradiction that at a local minimizer $\lambda_i, r_i >0$  for $i=1,2$ and hence by the complementary slackness property \cite[Proposition 3.3.5, (iv)]{bertsekas} the corresponding multipliers for the conditions $\lambda_1,\lambda_2 \geq 0$ and $0\leq r_1 \leq r_2 \leq R$ are $0$, since the constraints are inactive. Alternatively, we may consider the corresponding functions extended in a continuously differentiable way on the whole $\mathbb R^6$. 

The partial derivatives of \(u\) are

\[
\frac{\partial u}{\partial a}=1,
\qquad
\frac{\partial u}{\partial \kappa}=x, \quad
\frac{\partial u}{\partial \lambda_j}=-(x-r_j)_+
\text{ and for }
 x\neq r_j, \quad 
\frac{\partial u}{\partial r_j}
=
\lambda_j\1_{\{x>r_j\}}.
\]

Suppose first that \(f\) is a local minimizer with two changes of slope. Thus
\[
\lambda_1,\lambda_2>0
\qquad \text{and } \quad
0<r_1<r_2<R.
\]

% This means the parameter point is interior with respect to the kink constraints: no slope drop is zero, and no affine interval has zero length. The endpoint inequality \(\kappa^2-e^a\le0\) may or may not be active, and Fritz John handles either case.

By the Fritz--John necessary conditions, Lemma \ref{FJ}, there exist multipliers
$\tau\ge0,$ 
$\alpha,\beta\in\mathbb R,$ 
$\gamma\ge0,$
not all zero, such that the gradient of

\[
\tau h(f)
+
\alpha\left(\int_0^R f-1\right)
+
\beta\left(\int_0^R x f-1\right)
+
\gamma(\kappa-e^{\frac{a}2})
\]
vanishes. 

We first show that
$
\tau>0.$
Assume \(\tau=0\). Since \(\kappa-e^{\frac{a}{2}}\) does not depend on \(r_2\) or \(\lambda_2\), differentiating with respect to \(r_2\) and \(\lambda_2\) gives

\[
\int_{r_2}^R (\alpha+\beta x)f(x)\dd x=0
\quad \text{ and } \quad
\int_{r_2}^R (x-r_2)(\alpha+\beta x)f(x)\dd x=0.
\]
On \((r_2,R)\), the function \(\alpha+\beta x\) is affine. Writing 

\[
\alpha+\beta x=A+B(x-r_2),
\]
we have

\[
\int_{r_2}^R (\alpha+\beta x)^2f(x)\dd x
=
A\int_{r_2}^R (\alpha+\beta x)f(x)\dd x
+
B\int_{r_2}^R (x-r_2)(\alpha+\beta x)f(x)\dd x
=
0.
\]
Since \(f>0\) on \((r_2,R)\), it follows that
$ 
\alpha+\beta x=0
$ on $(r_2,R)$ and hence

\[
\alpha=\beta=0.
\]
Then differentiating with respect to \(a\) gives

\[
0=\gamma\frac{\partial}{\partial a}(\kappa-e^\frac{a}{2})
=
-\frac12\gamma e^{\frac{a}{2}}.
\]
Thus \(\gamma=0\), contradicting that the multipliers are not all zero. Therefore
$
\tau>0.
$
Dividing all multipliers by \(\tau\), we may assume \(\tau=1\).

Denote as shorthand

\[
W(x):=-(1+u(x))+\alpha+\beta x.
\]
Differentiating the Lagrangian with respect to \(\lambda_j\) gives

\begin{equation}
\int_{r_j}^R (x-r_j)W(x)f(x)\dd x=0.
\label{eq:danverm1}
\end{equation}

Differentiating with respect to \(r_j\) gives

\[
\lambda_j\int_{r_j}^R W(x)f(x)\dd x = 0.
\]
Since at \(x = r_j\) there is a change of slope, we must have \(\lambda_j>0\). Hence

\begin{equation}
\int_{r_j}^R W(x)f(x)\dd x=0.
\label{eq:danverm2}
\end{equation}

By \ref{eq:danverm1} and \ref{eq:danverm2} for $j=2$, we get

\[
\int_{r_2}^R Wf=0,
\qquad \text{and }
\int_{r_2}^R (x-r_2)Wf=0.
\]
On \((r_2,R)\), \(u\) is affine, so \(W\) is affine. Writing
\(
W(x)=A+B(x-r_2),
\)
we have

\[
\int_{r_2}^R W(x)^2f(x)\dd x
=
A\int_{r_2}^R Wf
+
B\int_{r_2}^R (x-r_2)Wf
=
0.
\]
Since \(f>0\) on \((r_2,R)\),
$
W=0$
on $(r_2,R).$
Thus
$W(r_2+)=0$ and also $W'_+(r_2)=0.$
But

\[
W'(x)=\beta-u'(x),
\]
which yields

\[
\beta -(\kappa -\lambda_1) = W'_-(r_2)=-\lambda_2.
\]
Therefore, on the middle interval \((r_1,r_2)\),
$W'(x) = -\lambda_2$ and so
\[
W(x)=W(r_2) - \int_{x}^{r_2}W'(t)dt = 0 - \int_{x}^{r_2}-\lambda_2dt = \lambda_2(r_2-x)>0.
\]
Also \(W=0\) on \((r_2,R)\). Hence

\[
\int_{r_1}^R W(x)f(x)\dd x
=
\int_{r_1}^{r_2} W(x)f(x)\dd x
>0.
\]
This contradicts \ref{eq:danverm2}
 with \(j=1\). Therefore a local minimizer cannot have two changes of slope.

Now suppose \(f\) is a local minimizer with exactly one change of slope:

\[
u(x)=a+\kappa x-\lambda(x-r)_+,
\qquad
\lambda>0,
\qquad
0<r<R.
\]
The same argument as in the previous case gives multipliers, with  nonzero multiplier for the entropy term, such that the gradient of the Lagrangian vanishes and hence a function

\[
W(x)=-(1+u(x))+\alpha+\beta x
\]
such that

\[
\int_r^R Wf=0,
\quad \text{and} \quad
\int_r^R (x-r)Wf=0.
\]
Since \(W\) is affine on \((r,R)\), the same square argument gives
$W=0$ on $(r,R).$
Thus
\[
W(r+)=0,
\quad \text{and} \quad
W'_+(r)=0.
\]
Differentiating $u$ we get
\[
W'_-(r)=-\lambda.
\]
Therefore, on \((0,r)\),

\[
W(x)=\lambda(r-x)>0.
\]
Now we differentiate the Lagrangian with respect to \(\kappa\). Since

\[
\frac{\partial u}{\partial \kappa}=x \quad \text{and} \quad \frac{\partial}{\partial\kappa}(\kappa-e^{\frac{a}2})=1,
\]
we get

\begin{equation} \label{secondcontr}
\int_0^R xW(x)f(x)\dd x+\gamma = 0.
\end{equation}
But \(W>0\) on \((0,r)\), \(W=0\) on \((r,R)\), and \(f>0\). Since also $\gamma\geq 0$

\[
\int_0^R xW(x)f(x)\dd x + \gamma
\geq
\int_0^r xW(x)f(x)\dd x
>0,
\]
contradicting \eqref{secondcontr}.

Thus a local minimizer cannot have one change of slope either.
Hence \(u\) is affine on its support, and

\[
f(x)=Ce^{\kappa x}\1_{(0,R)}(x),
\]
i.e., every local minimizer is one-piece log-affine. It remains to verify the desired inequality for such a density. Since $f$ is a mean-one density we must have 
\[
C=\frac{\kappa}{e^{\kappa R}-1},\quad \text{and} \quad
\frac{e^{\kappa R}(\kappa R-1)+1}{\kappa(e^{\kappa R}-1)} = 1.
\]
Writing \(t=\kappa R\),
we have
\[
h(f)=-\int_0^R f(x)\log f(x)\ dx
=
\log\Biggl(\frac{(e^t-1)^2}{e^t(t-1)+1}\Biggr)
-\frac{e^t(t-1)+1}{e^t-1}.
\]
The condition \((\log f)'(0)\le \sqrt{f(0)}\) is \(\kappa\le\sqrt C\), which is equivalent to $t\leq 1$. Differentiating the above expression for \(h\) with respect to \(t\), we get
\[
\frac{dh(f)}{dt}=
\frac{e^t(e^t t^2-(e^t-1)^2)}
{(e^t-1)^2(e^t(t-1)+1)}\le0,
\]
since \((e^t-1)^2\ge e^t t^2\). Therefore 
\[
h(f) \geq 2\log(e-1)-\frac1{e-1},
\]
which is the desired inequality. 



    The statement for centered densities follows immediately from Lemma \ref{lemmasatisfiescond}.
\end{proof}


%\section{Open Questions}
% Applying Theorem \ref{entropygrunbaumTh}, we get 
% $$
% h(X)\leq h(X^+) + \frac{H_2(p_+)}{p+}.
% $$
% Now noting that $H_2(x)$ is concave in $x \in [0,1]$, which implies that $\frac{H_2(x)}{x}$ is non-increasing in $x$, and applying Gr{\"u}nbaum's inequality, we obtain 
% $$
% h(X) \leq h(X^+) + e{H_2(\frac{1}{e})}.
% $$

% However, we believe this is not tight as our inequality is extremized for the negative exponential, while Gr{\"u}nbaum's inequality is extremized for the usual exponential. In view of the calculation \eqref{exponentialvalue}, we ask 
% \begin{question}
%     Let $X$ have a log-concave density $f$ with $\int_{\mathbb{R}}f(x)xdx = 0$. Is it true that 
%     $$
%     h(X^+) \geq h(X) - \frac{e}{e-1}H_2(\frac{1}{e})?
%     $$
%     With equality for positive exponent in the exponential. 
% \end{question}

%\begin{question}
%    Let $\Phi:\RL_+\to \RL$ be a continuous, convex function with $\Phi(0)=0.$ Is it true that 
  
 %   \[
    %\int\Phi(f(x))dx\le\int\Phi(f_{-}(x))dx
    %\]
%\end{question}


% \textcolor{blue}{$X$ could be negative no? i am confused...}
%%%%%%%%%%%%%%%%%%%%%%%%%%%%%%%%%%%%%

% {\color{blue} 
% \begin{theorem}
% Let \( X \) and \( Y \) be two independent random vectors in \( \mathbb{R}^n \). Then the following inequality holds
% \[
% I\left( P_n(X+Y) \right)^{-1} \geq I\left( P_n(X) \right)^{-1} + I\left( P_n(Y) \right)^{-1}
% \]
% where \( I(P_{n}(X)) \) corresponds to the Fisher information of the projection of the random vector onto the \( n \)-th coordinate, that is:
% \[
% I(P_n(X)) = \int_{\mathbb{R}^n} \frac{\langle \nabla f, e_n \rangle^2}{f(y)} \, dy,
% \]
% where \( f \) is the probability density of \( X \),  \( e_i \) is the standard basis vector in the \( i \)-th coordinate
% and \( \langle \cdot, \cdot \rangle \) denotes the usual inner product in \( \mathbb{R}^n \). We choose the \(n\)-th coordinate for simplicity; the theorem applies to any coordinate.
% \end{theorem}

% \begin{remark}
% Starting from the obtained inequality and the identity \( I(X) = \int_{\mathbb{S}^{n-1}} I(P_n(X)) \, d\sigma(u) \), it is possible to recover the Blachman-Stam inequality. Indeed,

% \[
% I(X+Y) = \int_{\mathbb{S}^{n-1}} I(P_n(X+Y)) \, d\sigma(u) \leq \int_{\mathbb{\mathbb{S}}^{n-1}} \left( I(P_n(X))^{-1} + I(P_n(Y))^{-1} \right)^{-1} \, d\sigma(u)
% \]

% \[
% \leq \left( \left( \int_{\mathbb{S}^{n-1}} I(P_n(X)) \, d\sigma(u) \right)^{-1} + \left( \int_{\mathbb{S}^{n-1}} I(P_n(Y))\, d\sigma(u) \right)^{-1} \right)^{-1}.
% \]

% %\textcolor{red}{Why the identity $I(X)=\int_{S^{n-1}} I(P_{n}(X))d\sigma(u)$ is true ?}
% \end{remark}

% %\begin{remark}
% %Given the analogies between convex geometry and information theory, the Fisher information $I(X)$ is the analog of $\frac{\lvert \partial A \rvert}{\lvert A \rvert}$ for a compact set, the obtained entropy inequality is the analog of  the following one 
% %\[
% %\frac{\lvert P_n(A + B) \rvert}{\lvert \partial(P_n(A + B)) \rvert} \geq \frac{\lvert A \rvert}{\lvert \partial A \rvert} + \frac{\lvert B \rvert}{\lvert \partial B \rvert},
% %\]
% %where $A$ and $B$ are compact sets in $\mathbb{R}^n$ and $P_n$ is the projection onto  $e_n$.
% %\textcolor{red}{Not completely sure... but seems true , a little bit vague. In any case, I think we were told by Matthieu that this last conjecture was false}


% %\end{remark}


% \begin{proof}

% Let $X$ be a random vector as stated. Let $T_m: \mathbb{R}^n \to \mathbb{R}^n$ be an operator defined as follows

% \[
% T_m = I_n + \left( \frac{1}{m} - 1 \right) e_n e_n^T .
% \]
% Then, given that $X \sim f$,
% we know that 
% \[
% T_{m}(X)\sim \frac{1}{\lvert T \rvert} f(T_{m}^{-1} x) := f_{T_{m}}.
% \]
% Thus, the Fisher information $I(T_{m}(X))$ is given by
% \[
% I(T_{m}(X)) = \int_{\mathbb{R}^{n}} \frac{\lvert \nabla f_{T_{m}}(x) \rvert^2}{f_{T_{m}(x)}} \, dy
% = \int_{\mathbb{R}^{n}} \frac{\lvert (T^{-1})^T \nabla f(y) \rvert^2}{f(y)} \, dy
% = \int_{\mathbb{R}^{n}} \frac{\sum_{i=1}^{n-1} (\partial_i f)^2 + m^2 (\partial_n f)^2}{f(y)} \, dy.
% \]

% Therefore, 
% \[
% \lim_{m \to \infty} \frac{I(T_{m}(X))}{m^2} = \int_{\mathbb{R}^{n}} \frac{(\partial_n f)^2}{f(y)} \, dy = I(P_n(X)).
% \]

% Now, applying the Blachman-Stam inequality to the independent random variables $T_m(X)$ and $T_m(Y)$ , multiplying by $m^2$ and letting $m$ tend to infinity, the result follows.
% \end{proof}
%  }
 
 
 

 \newpage 
\section*{Appendix}

\subsection*{Basic information theoretic properties}

In this appendix we summarize a few properties of entropy and mutual information. 

We recall here that the differential entropy of a vector $X$ with density $f$ on $\mathbb{R}^n$ is 
$$
h(X) = h(f) := -\int_{\mathbb{R}^n}{f(x)\log{f(x)} dx} \quad \in [-\infty, \infty].
$$
A change of variables in the integral shows that the differential entropy satisfies, for any invertible matrix $A \in \mathbb{R}^{n \times n}$, the scaling property
\begin{equation}
    h(AX) = h(X) + \log{|\det{(A)}|}.
\end{equation}

If $X \in \mathbb{R}^n,Y \in \mathbb{R}^m$ are two random variables (resp. vectors) with joint density $f_{X,Y}$ on $\mathbb{R}^{m+n}$ then the conditional density of $X$ given $Y$ exists a.s. and is given for any $y \in \mathbb{R}^m$ by 
$$
f_{X|Y}(x|y) := \frac{f_{X,Y}(x,y)}{f_Y(y)},
$$
where $f_Y(y) = \int{f_{X,Y}(x,y) dx}$ is the marginal of $Y$.
Then the joint and conditional entropies are defined respectively as 
$$
h(X,Y) := -\int_{\RL^{m+n}}{f(x,y)\log{f(x,y)}dxdy}
$$
and
$$
h(X|Y) := -\int_{\RL^{n+m}}{f_{X,Y}(x,y)\log{f_{X|Y}(x|y)}dxdy}.
$$
The conditional entropy can also be re-written as 
\begin{align} \nonumber
h(X|Y) &= -\int_{\RL^{n+m}}{f_{X,Y}(x,y)\log{f_{X|Y}(x|y)}dxdy} \\ \label{condHexpression}
&= -\int_{\RL^{m}} {f_Y(y)\int_{\RL^n}{f_{X|Y}(x|y)\log{f_{X|Y}(x|y)}dx}dy} \\ \label{condHaverage}
&= \int_{\RL^m}{f_Y(y)h(X|Y=y)}dy, 
\end{align}
where we denote, for $y \in \RL^m$, $h(X|Y=y) := h\bigl(f_{X|Y}(\cdot|y)\bigr).$ More generally, if $Y$ does not have a density, \eqref{condHaverage} can be replaced by 
\begin{equation} 
h(X|Y) = \int{h(X|Y=y) d\mu_Y(y)}
\end{equation}
where $\mu_Y$ stands for the law of $Y$.

The chain rule for entropy \cite{cover:book} is 
\begin{equation} \label{chainrule}
h(X,Y) = h(Y) + h(X|Y) = h(X) + h(Y|X).
\end{equation}

If $(X,Y)$ has joint law $\mu_{X,Y}$ and $X,Y$ have marginal laws $\mu_X$ and $\mu_Y$ respectively, the {\em mutual information} between $X$ and $Y$ is 
\begin{equation} \label{mutualinfodef}
I(X;Y) := \int d\mu_{X,Y}\log{\frac{d\mu_{X,Y}}{d\mu_X \times d\mu_Y}} \geq 0.
\end{equation}
When $(X,Y)$ has a joint density
$$
I(X;Y) = h(X) - h(X|Y),
$$
while if $Y$ is discrete and $X$ has conditional density given $Y$ then 
$$
I(X;Y) = h(X) - \sum_yp_Y(y)h(X|Y=y).
$$
The positivity of mutual information gives the standard property 
\begin{equation} \label{conditionreduceentropy}
    h(X|Y) \leq h(X).
\end{equation}







\bibliographystyle{abbrv}
\bibliography{references}

\end{document}
